\section{Higher factorization categories and chiral derived centres at \texorpdfstring{$d\geq 2$}{d>=2}}\label{app:higher-factorization-categories}

The local theorem (Theorem~\ref{thm:main-local}) is the
formal-Darboux fibre of a higher open--closed problem.  At the chosen
Dirac brane vacuum \(b_p\in\mathcal{C}^{\mathrm{op}}_\partial\), the
Hamiltonian BF bulk gives Chevalley--Eilenberg cochains, the boundary
brane gives the \(E_1\)-Hochschild complex of
\(A^{\mathrm{cl}}_{\partial,N}\), and the trace map
\(J(f)=\operatorname{Tr}f(\phi_1,\phi_2)\) gives the coordinate
generator assignment whose centre lift is the closed-to-open morphism
problem.  The Mixed Holomorphic--Topological Deligne conjecture asserts
that, before choosing the vacuum, this local Hochschild chart is the
restriction of the chiral derived centre of the primitive open boundary
category at \(d=2\).

The higher factorization-categorical setting consists of holomorphic
factorization algebras, mixed holomorphic--topological factorization
algebras, chiral derived centre objects for \(d\geq2\), and the
formal-Darboux restriction functor.  At \(d=1\) this language reduces
to Beilinson--Drinfeld chiral algebras on curves
\cite{beilinson-drinfeld-chiral}.  At \(d=2\), the formal-Darboux
trace theorem gives the trace-sector stalk; the formal polydisk also
carries the all-arity tree-kernel current operad constructed in
Theorem~\ref{app:thm:formal-polydisk-all-arity-current-operad}, whose
binary generator is the diagonal local-cohomology singular product of
Theorem~\ref{thm:formal-two-dimensional-holomorphic-ope}.  The global
chiral \(E_2\)-Deligne problem is the promotion of this formal current
operad to a Morita chiral \(E_2\)-centre Maurer--Cartan lift with
Ran/Weiss descent and QME compatibility.

\subsection{\texorpdfstring{Factorization $E_d$-algebras on smooth complex $d$-folds}{Factorization Ed-algebras on smooth complex d-folds}}\label{app:hfc:fact-Ed}

Let $X$ be a smooth complex variety of complex dimension $d$.  Write
$\mathrm{Ran}(X)$ for the Ran space of finite subsets of $X$.  Let
$\mathrm{Mfd}^{\mathrm{lc}}(X)$ be the symmetric monoidal $\infty$-category
whose objects are open subsets $U\subset X$ and whose monoidal product
is disjoint union of subsets with disjoint closures.

\begin{defn}[Holomorphic factorization algebra]\label{app:defn:hol-fact-Ed}
A \emph{holomorphic factorization algebra} on \(X\) is a
Costello--Gwilliam prefactorization algebra valued in cochain
complexes, with holomorphic locality data and Weiss descent.  It
assigns a cochain complex \(U\mapsto\mathcal F(U)\) to each open
\(U\subset X\), structure maps
\[
  \mathcal F(U_1)\otimes\cdots\otimes\mathcal F(U_n)
   \longrightarrow
   \mathcal F(V)
\]
for pairwise disjoint opens \(U_i\subset V\), and extension maps for
inclusions, satisfying the factorization axioms of Costello--Gwilliam
\cite{costello-gwilliam}*{Vol.~I, Defn.~3.1.4}.  The holomorphic
condition means that the local cochain models are Dolbeault
resolutions and the structure maps commute with the holomorphic
differential.  An \(E_d\)-algebra object in the category of such
factorization algebras is extra structure.
\end{defn}

\begin{rmk}[Comparison presentations]
A holomorphic factorization algebra on $X$ has three comparison
languages
\cite{beilinson-drinfeld-chiral}\cite{francis-thesis}\cite{lurie-ha}*{\S 5.5}:
\begin{enumerate}
  \item \emph{Costello--Gwilliam.}  A factorization algebra in the sense
    of \cite{costello-gwilliam}*{Vol.~I, \S 3}, with the holomorphic
    locality differential.
  \item \emph{Factorisable $\mathcal D$-module on
    $\mathrm{Ran}(X)$.}  A factorisable $\mathcal D$-module
    $\mathcal A$ on $\mathrm{Ran}(X)$.  At
    $d=1$ this is the algebro-geometric construction of chiral
    algebras on curves \cite{beilinson-drinfeld-chiral}*{\S 3}.  At
    $d\geq2$ the corresponding Ran presentation is an additional
    comparison datum.
  \item \emph{Lurie--Ayala--Francis.}  An $E_d$-algebra object of the
    $\infty$-category of factorisation algebras on $X$ in the sense of
    \cites{ayala-francis}*{}{lurie-ha}*{\S 5.5.4}.
\end{enumerate}
At $d=1$ these presentations reduce to chiral algebras on curves in the
algebro-geometric sense of \cite{beilinson-drinfeld-chiral}*{\S 3}.
For $d\geq 2$ the Costello--Gwilliam holomorphic
factorization-algebra presentation is the definition of
the closed-string sector classical observables.  The Ran
\(\mathcal D\)-module comparison and the Lurie--Ayala--Francis
comparison are hypotheses in the global chiral \(E_d\)-Deligne
problem isolated in
\S\ref{ssec:open-chiral-Ed-Deligne}.
\end{rmk}

\begin{defn}[Category of holomorphic factorization algebras]\label{app:defn:HolFA-d-cat}
Let $\mathrm{HolFA}_d(X)$ denote the symmetric monoidal $\infty$-category
whose objects are holomorphic factorization algebras on $X$ in the
sense of Definition~\ref{app:defn:hol-fact-Ed}, and whose morphisms are
natural transformations of prefactorization diagrams that intertwine
the factorization structure maps.  Its monoidal product is the
pointwise tensor product of cochain complexes.  An \(E_d\)-algebra
object of \(\mathrm{HolFA}_d(X)\) is extra structure.
\end{defn}

\subsection{Mixed holomorphic-topological refinement}\label{app:hfc:mixed-HT}

The local geometry is
$X=\R^k_{\mathrm{top}}\times M_{\mathrm{hol}}$ with $M$ a smooth
complex variety of complex dimension $d$ ($k=2$, $d=2$ in the body).
The factorization algebra carries three independent locality
differentials: support locality on $\R^k$, holomorphic locality on
$M$, and topological / locally constant in $\R^k$.

\begin{defn}[Mixed-HT factorization $E_d^{\mathrm{HT}}$-algebra]\label{app:defn:mixed-HT-fact}
A \emph{mixed-HT factorization $E_d^{\mathrm{HT}}$-algebra} on
$X=\R^k_{\mathrm{top}}\times M_{\mathrm{hol}}$ is a
Costello--Gwilliam prefactorization algebra on \(X\), valued in
cochain complexes, satisfying Weiss descent and carrying the three
locality data:
\begin{enumerate}[label=(\roman*)]
  \item \emph{Support locality on $\R^k$.}  For any compact embedding
    $K\subset U\subset\R^k$, the section $\mathcal F$ is supported on
    $K$ when restricted to $U\setminus K$ extends to zero.
  \item \emph{Topological locality on $\R^k$.}  For any open
    $U\subset\R^k$ and any inclusion $U_0\subset U$ with
    $\pi_0(U_0)\to\pi_0(U)$ a bijection, the structure map
    $\mathcal F(U_0\times B)\to\mathcal F(U\times B)$ is a
    quasi-isomorphism for every $B\subset M$.
  \item \emph{Holomorphic locality on $M$.}  Restriction to the
    holomorphic factor is a holomorphic factorization algebra
    in the sense of Definition~\ref{app:defn:hol-fact-Ed}.
\end{enumerate}
An \(E_d^{\mathrm{HT}}\)-operadic action on such a factorization
algebra is additional structure.
\end{defn}

\begin{defn}[Category $\mathrm{ChirAlg}^{\mathrm{HT}}_{k,d}(X)$]\label{app:defn:ChirAlg-HT}
Let $\mathrm{ChirAlg}^{\mathrm{HT}}_{k,d}(X)$ denote the symmetric
monoidal $\infty$-category whose objects are mixed-HT factorization
$E_d^{\mathrm{HT}}$-algebras on $X$ in the sense of
Definition~\ref{app:defn:mixed-HT-fact}, and whose morphisms intertwine
all three locality structures.  At $k=0$ this reduces to
$\mathrm{HolFA}_d(M)$; at $d=0$ it reduces to topological factorization
algebras on $\R^k$.  The formal-Darboux local geometry is $k=2$, $d=2$.
\end{defn}

\subsection{The chiral derived centre functor}\label{app:hfc:chiral-Z-der}

\begin{defn}[Chiral derived centre at $d=1$]\label{app:defn:Z-der-ch-d1}
For a chiral algebra $\mathcal A\in\mathrm{HolFA}_1(C)$ on a smooth
curve $C$, the \emph{chiral derived centre}
$Z^{\mathrm{der}}_{\mathrm{ch}}(\mathcal A)$ is the chiral Hochschild
cochain complex
\[
  Z^{\mathrm{der}}_{\mathrm{ch}}(\mathcal A)
   =C^\bullet_{\mathrm{ch}}(\mathcal A,\mathcal A)
   :=\mathrm{RHom}_{\mathcal A\otimes\mathcal A^{\mathrm{op}}}(\mathcal A,\mathcal A)
\]
in the chiral monoidal structure of
\cite{beilinson-drinfeld-chiral}*{\S 3.4}.  It carries a natural
$E_2$-algebra structure (chiral Deligne theorem at \(d=1\);
\cite{beilinson-drinfeld-chiral}*{\S 4}, Lurie HA
\cite{lurie-ha}*{Thm.~5.3.1.14}, Costello--Gwilliam
\cite{costello-gwilliam}*{Vol.~I, \S 6}).  This is the proved $d=1$
chiral Deligne theorem of Vol I.
\end{defn}

\begin{defn}[Chiral derived-centre model at $d\geq 2$]\label{app:defn:Z-der-ch-dgeq2}
For a holomorphic factorization algebra
$\mathcal A\in\mathrm{HolFA}_d(X)$ on a smooth complex variety $X$ of
complex dimension \(d\geq2\), the \emph{chiral derived-centre model
at \(d\)} is the endomorphism object
\[
  Z^{\mathrm{der}}_{E_d^{\mathrm{ch}}}(\mathcal A)
   :=
   \mathrm{RHom}_{\mathcal A\otimes\mathcal A^{\mathrm{op}}}(\mathcal A,\mathcal A)
\]
in a chosen \(\infty\)-category of factorization \(E_d\)-algebra
bimodules on \(X\).  The notation records the chiral centre;
the higher chiral monoidal structure and its \(E_{d+1}^{\mathrm{ch}}\)
operations are extra data, not part of this definition.  At \(d=1\)
this construction reduces to Definition~\ref{app:defn:Z-der-ch-d1}.

\emph{The obstruction at $d\geq 2$.}  The existence of an
$E_{d+1}^{\mathrm{ch}}$-structure on
$Z^{\mathrm{der}}_{E_d^{\mathrm{ch}}}(\mathcal A)$ extending the
$E_2$-structure of the $d=1$ case is precisely the datum asserted by
the chiral $E_d$-Deligne conjecture
(\S\ref{ssec:open-chiral-Ed-Deligne}; Lurie
\cite{lurie-ha}*{\S 5.3.1}); partial constructions are due to
Hennion (chiral Lie algebroids on smooth schemes) and Francis
(factorisable co-operads).  The Mixed Holomorphic--Topological Deligne
conjecture at $d=2$ is the corresponding statement for
$\mathrm{ChirAlg}^{\mathrm{HT}}_{2,2}(\R^2\times\C^2)$.
\end{defn}

\begin{defn}[Morita mixed-HT chiral derived centre]\label{app:defn:Z-der-HT}
Let \(X=\R^k_{\mathrm{top}}\times \C^d_{\mathrm{hol}}\), and let
\(\mathcal C\) be a mixed-HT open boundary category.  Fix the
Tate/pro-Matlis complete admissible enlargement
\(\widehat{\mathcal C}_{\mathrm{adm}}\) used for boundary observables.
Let
\[
  \mathrm{Bimod}^{\mathrm{HT}}_{\operatorname{Ran}(X)}
       (\mathcal C,\mathcal C)
\]
be the stable \(\infty\)-category of complete factorisation bimodules on
\(\operatorname{Ran}(X)\) with commuting left and right
\(\widehat{\mathcal C}_{\mathrm{adm}}\)-actions, continuous for the
Tate/pro-Matlis topology and satisfying Weiss descent on the brane-link
boundary.  Let \(\Delta_{\mathcal C}\) denote the diagonal
\((\mathcal C,\mathcal C)\)-bimodule, equivalently the identity module
category of \(\widehat{\mathcal C}_{\mathrm{adm}}\).  The
\emph{Morita mixed-HT chiral derived centre} is
\[
  Z^{\mathrm{der},\mathrm{Mor}}_{E_d^{\mathrm{HT}}}(\mathcal C)
  :=
  \operatorname{RHom}_{
    \mathrm{Bimod}^{\mathrm{HT}}_{\operatorname{Ran}(X)}
      (\mathcal C,\mathcal C)}
  (\Delta_{\mathcal C},\Delta_{\mathcal C}).
\]
When no ambiguity is possible, the manuscript writes this Morita object
as \(Z^{\mathrm{der}}_{E_d^{\mathrm{HT}}}(\mathcal C)\).
The notation records the centre object.  For \(d\geq2\), a mixed-HT
chiral \(E_{d+1}\)-algebra structure on this object is an additional
chiral \(E_d\)-Deligne datum.

If \(b\in\mathcal C\) is a dualisable boundary object and
\(A_b=\operatorname{End}_{\mathcal C}(b)\), restriction of bimodules to
the full subcategory generated by \(b\) gives the Hochschild chart
\[
  \operatorname{ev}_b\colon
  Z^{\mathrm{der},\mathrm{Mor}}_{E_d^{\mathrm{HT}}}(\mathcal C)
  \longrightarrow
  C^\bullet_{\mathrm{ch},\mathrm{HT}}(A_b,A_b).
\]
This chart is an equivalence only under a separate Morita-density
hypothesis on \(b\).  At the \(N\)-Dirac-brane formal-Darboux vacuum
\(b_p\), the theorem computes the trace-sector image of
\(\operatorname{ev}_{b_p}\), not the whole category-level centre.
\end{defn}

\subsection{The formal-Darboux stalk functor}\label{app:hfc:darboux-stalk-functor}

For $X=\R^k_{\mathrm{top}}\times\C^d_{\mathrm{hol}}$ and a closed point
$p\in\C^d$, the formal-Darboux stalk operation is restriction to the
formal completion at $p$.

\begin{defn}[Formal-Darboux stalk functor]\label{app:defn:Darboux-stalk-functor}
The \emph{formal-Darboux stalk functor} at $p$ is the symmetric monoidal
$\infty$-functor
\[
  (-)|_{p}\colon
  \mathrm{ChirAlg}^{\mathrm{HT}}_{k,d}(\R^k\times\C^d)
   \longrightarrow
  \mathrm{ChirAlg}^{\mathrm{HT}}_{k,d}(\R^k\times\widehat{\C^d}_p),
\]
sending $\mathcal F\mapsto\mathcal F|_{\R^k\times\widehat{\C^d}_p}$, the
restriction of the factorisable structure to the formal neighbourhood
$\widehat{\C^d}_p=\operatorname{Spf}\widehat{\mathcal O}_{\C^d,p}$.  The
holomorphic Darboux theorem trivialises the germ of $\C^d$ at $p$ to
the formal symplectic polydisk in coordinates $(z_1,\ldots,z_d)$, so
the value at $p$ depends on $p$ only through its formal isomorphism
class --- equivalently, only through the formal Darboux coordinates.
\end{defn}

\begin{prop}[Equivariance under formal symplectomorphism]\label{app:prop:Darboux-equivariance}
The formal-Darboux stalk functor of
Definition~\ref{app:defn:Darboux-stalk-functor} is equivariant under the
formal symplectomorphism group $\mathrm{Symp}^{\mathrm{form}}(\C^d)$
acting on the source category by formal coordinate change at $p$ and
on the target by the same change of formal Darboux coordinates.  At
$d=2$ on the holomorphic-symplectic factor, the stalk is canonical up
to formal Darboux change.
\end{prop}

\begin{proof}
The Darboux theorem in formal coordinates trivialises any holomorphic
symplectic form to $\omega_p=\sum_i dz_{2i-1}\wedge dz_{2i}$, with the
trivialisation defined uniquely up to the formal symplectomorphism
group.  The factorisation structure pulls back along this
trivialisation by symmetric monoidal naturality.
\end{proof}

\subsection{The bulk closed-string sector as an
\texorpdfstring{$E_2^{\mathrm{HT}}$-factorization}{E2-HT-factorization}
algebra}\label{app:hfc:bulk-as-fact-cat}

\begin{constr}[Bulk classical observables on $\R^2\times\C^2$]\label{app:constr:bulk-fact-Ed}
Let $X=\R^2_{\mathrm{top}}\times\C^2_{\mathrm{hol}}$ and
$\mathfrak g=\mathfrak h\ltimes\mathfrak h^\vee_{\mathrm{cont}}[1]$ as
in \S\ref{ssec:hamiltonian-bf-algebra}.  The closed-string sector
classical observables factorization algebra
$\mathcal A^{\mathrm{cl}}_{\mathrm{bulk}}\in
\mathrm{ChirAlg}^{\mathrm{HT}}_{2,2}(X)$ is the assignment
\[
  U\times B\longmapsto
   C^\bullet_{\mathrm{CE},\mathrm{cont}}\bigl(
     \Omega^\bullet_c(U)\widehat\otimes
     \Omega^{0,\bullet}_c(B)\widehat\otimes\mathfrak g\bigr),
\]
for $U\subset\R^2$ and $B\subset\C^2$ open, with $E_2$-algebra structure
from the symmetric algebra structure of the continuous CE cochains and
factorisation isomorphisms inherited from extension-by-zero on
disjoint opens.  The three locality differentials of
Definition~\ref{app:defn:mixed-HT-fact} are the de Rham differential
on $U$, the Dolbeault differential $\bar\partial$ on $B$, and the
internal CE differential dual to the $\mathfrak g$-bracket.
\end{constr}

\begin{prop}[Formal-Darboux stalk of the bulk]\label{app:prop:bulk-stalk}
At a closed point $p\in\C^2$, the formal-Darboux stalk of
$\mathcal A^{\mathrm{cl}}_{\mathrm{bulk}}$ is the continuous CE cochain
algebra of the shifted-cotangent BF Lie algebra at the formal polydisk:
\[
  \mathcal A^{\mathrm{cl}}_{\mathrm{bulk}}|_p
   =C^\bullet_{\mathrm{CE},\mathrm{cont}}(\mathfrak g),
   \qquad
   \mathfrak g=\mathfrak h\ltimes\mathfrak h^\vee_{\mathrm{cont}}[1],
   \quad
   \mathfrak h=\C[[z_1,z_2]]/\C\cdot 1.
\]
The formal coordinate algebra
\(\widehat{\mathcal O}_{\C^2,p}\) gives the Hamiltonian labels
\(\mathfrak h=\widehat{\mathcal O}_{\C^2,p}/\C\cdot1\); it is not the
closed CE cochain algebra and it is not the brane endomorphism
algebra.  The arity-two singular operation on the closed stalk is the
native diagonal singular product
\[
  \mathsf{J}_x(z)\mathsf{J}_y(w)
  \sim
  K^{(2)}_\Delta(z,w)\mathsf{J}_{[x,y]_{\mathfrak g}}(w),
  \qquad x,y\in\mathfrak g,
\]
with
\[
  K^{(2)}_\Delta=\frac{d\zeta_1d\zeta_2}{\zeta_1\zeta_2},\qquad
  \zeta_i=z_i-w_i .
\]
\end{prop}

\begin{proof}
Apply Definition~\ref{app:defn:Darboux-stalk-functor} to
Construction~\ref{app:constr:bulk-fact-Ed}: the de Rham and Dolbeault
Poincar\'e contractions on the formal polydisk reduce the compactly
supported sheaf data to the continuous CE cochains of $\mathfrak g$ at
$p$.  Extension by zero gives the regular product on disjoint formal
polydisks.  When two formal insertions collide, the Dolbeault
resolution of the diagonal leaves the local-cohomology quotient
\[
  H^2_\Delta(\C[[z_1,z_2,w_1,w_2]])\,d\zeta_1d\zeta_2,
\]
and the CE bracket of $\mathfrak g$ gives exactly the singular product
of Theorem~\ref{thm:formal-two-dimensional-holomorphic-ope}.  Thus the
formal completion retains the diagonal collision operation even though
its disjoint-open product is the extension-by-zero product.
\end{proof}

\subsection{The brane open category as factorization
\texorpdfstring{$E_1^{\mathrm{HT}}$}{E1-HT}-module}\label{app:hfc:brane-as-fact-cat}

\begin{defn}[Boundary open category]\label{app:defn:brane-cat}
The \emph{boundary open category} of the local theory is
\(\mathcal{C}^{\mathrm{op}}_\partial\), the category of boundary
conditions on the brane-link boundary $\partial X$ of the Dirac brane
defect
$L=\R_t\times\{0_s\}\times\{p\}
\subset X=\R^2_{\mathrm{top}}\times\C^2_{\mathrm{hol}}$
--- the link sphere bundle over the worldline $\R_t$, equivalently the
boundary of an excised tubular neighbourhood
$X\setminus\nu(L)$.  The bulk factorization algebra acts on
\(\mathcal{C}^{\mathrm{op}}_\partial\) by the bulk--boundary OPE.  A
module-category model for this action is extra structure, not the
definition of the open category.
\end{defn}

\begin{constr}[Formal-Darboux vacuum and brane $E_1$-algebra]\label{app:constr:brane-vacuum-E1}
The \emph{formal-Darboux vacuum} at the closed point $p\in\C^2$ is
an object \(b_p\in\mathcal{C}^{\mathrm{op}}_\partial\) equipped with
the formal coordinate chart
\(\widehat{\mathcal O}_{\C^2,p}\cong\C[[z_1,z_2]]\).  The chart gives
the substitution \(z_i\mapsto\phi_i\).  The derived endomorphism
algebra of a skyscraper brane is the corresponding Ext--Koszul
algebra, not the formal function ring itself.  At $N$ Dirac branes
over $b_p$ the brane $E_1$-algebra of classical observables on the
worldline $\R_t$ is
\[
  A^{\mathrm{cl}}_{\partial,N}
   =C^\bullet_{\mathrm{CE}}\bigl(\mathfrak{gl}_N,\,
       \mathrm{Sym}(\mathfrak{gl}_N\oplus\mathfrak{gl}_N
                    \oplus\mathfrak{gl}_N[1])\bigr),
   \qquad
   Q\psi=[\phi_1,\phi_2],
\]
with derived Chan--Paton substitution \(z_i\mapsto\phi_i\) and
Koszul differential \(Q\psi=[\phi_1,\phi_2]\).  On \(H^0\) this
substitution factors through the commuting-pair locus
(Definition~\ref{def:formal-darboux-stalk-chart}).  Curve chiral
algebras such as affine Kac--Moody, \(\beta\gamma\), Virasoro, and
\(W\)-algebras enter through a curve restriction and anomaly
matching before comparison with this formal disk.
\end{constr}

\begin{prop}[Brane $E_1$-Hochschild as $E_2$-algebra]\label{app:prop:brane-HH-E2}
The derived $E_1$-centre
$Z^{\mathrm{der}}_{E_1}(A^{\mathrm{cl}}_{\partial,N})$ is naturally an
$E_2$-algebra, equivalent to the Hochschild cochain complex
$C^\bullet_{\mathrm{Hoch}}(A^{\mathrm{cl}}_{\partial,N},A^{\mathrm{cl}}_{\partial,N})$;
its cohomology is
$HH^\bullet(A^{\mathrm{cl}}_{\partial,N},A^{\mathrm{cl}}_{\partial,N})$
with Gerstenhaber bracket and cup product.  This is Deligne's
$E_1$-Hochschild theorem
(Tamarkin~\cite{tamarkin-deligne},
McClure--Smith~\cite{mcclure-smith},
Kontsevich--Soibelman~\cite{kontsevich-soibelman-deformation},
Lurie~\cite{lurie-ha}*{Thm.~5.3.1.14}).
\end{prop}

\subsection{The Mixed Holomorphic--Topological Deligne conjecture as a
functorial statement}\label{app:hfc:mixed-HT-Deligne}

\begin{conj}[Mixed Holomorphic--Topological Deligne, functorial form]\label{app:conj:mixed-HT-Deligne}
For the bulk closed-string sector
$\mathcal A^{\mathrm{cl}}_{\mathrm{bulk}}\in\mathrm{ChirAlg}^{\mathrm{HT}}_{2,2}(\R^2\times\C^2)$
of Construction~\ref{app:constr:bulk-fact-Ed} and the boundary open
category $\mathcal{C}^{\mathrm{op}}_\partial$ of
Definition~\ref{app:defn:brane-cat}, there is a canonical equivalence
  \[
	  \mathcal A^{\mathrm{cl}}_{\mathrm{bulk}}
	   \;\simeq\;
	   Z^{\mathrm{der},\mathrm{Mor}}_{E_2^{\mathrm{HT}}}
	   \bigl(\mathcal{C}^{\mathrm{op}}_\partial\bigr)
	\]
of mixed-HT chiral algebras at $d=2$, where the right-hand side is the
Morita mixed-HT chiral derived centre of
Definition~\ref{app:defn:Z-der-HT}.  After choosing a dualisable
\(b\in\mathcal C^{\mathrm{op}}_\partial\), the evaluation
\(\operatorname{ev}_b\) gives the Hochschild chart
\(C^\bullet_{\mathrm{ch},\mathrm{HT}}(A_b,A_b)\); the chart is not a
definition of the category-level centre.

\emph{Known cases.}
\begin{itemize}
  \item At $d=1$ (replacing $\C^2$ by a smooth curve $C$): proved
    (chiral Deligne theorem at \(d=1\),
    Vol I~\cite{lorgat-vol1}; \cite{beilinson-drinfeld-chiral}).
  \item At $d\geq 2$ the datum is the chiral $E_d$-Deligne conjecture in
    the sense of Lurie~\cite{lurie-ha}*{\S 5.3.1}.  Its mixed-HT form is
    the holomorphic-topological enhancement.  Partial constructions are
    recorded in Hennion's chiral Lie algebroids on smooth schemes and
    Francis's factorisable co-operads.
	  \item \emph{Formal-Darboux stalk:} the trace-sector
	    map
	    \[
	      \Phi_N^{\mathrm{gen}}\colon
	      C^\bullet_{\mathrm{CE},\mathrm{cont}}(\mathfrak g)
	      \dashrightarrow
	      C^\bullet_{\mathrm{Hoch}}\!
      (A^{\mathrm{cl}}_{\partial,N},A^{\mathrm{cl}}_{\partial,N}),
      \qquad
      c_f\mapsto\theta_f,\quad u_f\mapsto J(f),
	    \]
	    at \(N\) Dirac branes over the formal-Darboux vacuum \(b_p\).
	    On the stable scalar-reduced trace algebra its coordinate shadow
	    \(\Phi_{\mathrm{coord}}^{\mathrm{st}}\) is the
	    Chevalley--Eilenberg/polyvector coordinate isomorphism after the
	    admissible Tate/pro-Matlis completion.  The whole finite-\(N\)
    derived-centre equivalence is governed by Hochschild centrality
    homotopies, finite-trace relations, and descent/QME hypotheses
    recorded below.  The
    formal-coordinate arity-two singular product is computed by
    Theorem~\ref{thm:formal-two-dimensional-holomorphic-ope}.
\end{itemize}
\end{conj}

\begin{thm}[Formal-Darboux trace-sector stalk of Conjecture~\ref{app:conj:mixed-HT-Deligne}]\label{app:thm:stalk-of-conj}
At $N$ Dirac branes over the formal-Darboux vacuum
$b_p\in\mathcal{C}^{\mathrm{op}}_\partial$, the formal-Darboux stalk
of Conjecture~\ref{app:conj:mixed-HT-Deligne} contains the
trace-sector morphism
\[
  \mathcal A^{\mathrm{cl}}_{\mathrm{bulk}}|_p
   \simeq
	   C^\bullet_{\mathrm{CE},\mathrm{cont}}(\mathfrak g)
	   \dashrightarrow_{\Phi_N^{\mathrm{gen}}}
		   C^\bullet_{\mathrm{Hoch}}\bigl(A^{\mathrm{cl}}_{\partial,N},A^{\mathrm{cl}}_{\partial,N}\bigr)
\]
with formal Darboux coordinate expression
\(J(f)=\overline{\operatorname{Tr}\,f(\phi_1,\phi_2)}\) on polynomials,
extended to
\(\mathfrak h=\widehat{\mathcal O}_{\C^2,p}/\C\cdot 1\) only after
formal completion, and reduced closed-to-open dictionary
\(c_f\mapsto\theta_f\in C^1\), \(u_f\mapsto J(f)\cdot(-)\in C^0\) on
Chevalley--Eilenberg generators of
$\mathfrak g=\mathfrak h\ltimes\mathfrak h^\vee_{\mathrm{cont}}[1]$.
After scalar reduction and passage to the stable trace algebra the
coordinate shadow \(\Phi_{\mathrm{coord}}^{\mathrm{st}}\) is the
admissible CE/PV coordinate isomorphism.  For fixed
finite \(N\) it is not an equivalence onto the whole Hochschild
complex.
The arity-two formal diagonal operation at the stalk is the singular
product of
Theorem~\ref{thm:formal-two-dimensional-holomorphic-ope}.
The chain-level lift \(\Phi_N^{\mathrm{Hoch}}\) to a strict
\(P_0\)-factorization-centre
morphism exists under compatible null-homotopies of
$\operatorname{Ob}_{\mathrm{cent}/\mathrm{QME}/\mathrm{desc}}$
(Theorem~\ref{thm:app-unreduced-cotangent-lift-obstruction}).
\end{thm}

\begin{proof}
Combine Proposition~\ref{app:prop:bulk-stalk} (formal-Darboux stalk of
the bulk),
Proposition~\ref{app:prop:brane-HH-E2} (Deligne's $E_1$-Hochschild
theorem for the brane $E_1$-algebra), and the trace-sector
closed-to-open construction of Theorem~\ref{thm:main-local}.  The
trace-map computation is
Theorem~\ref{thm:formal-local-universal-property} (Procesi--Razmyslov +
Loday--Quillen--Tsygan); the dictionary is
Theorem~\ref{thm:reduced-ce-pv-central-operation}; the arity-two
diagonal operation is
Theorem~\ref{thm:formal-two-dimensional-holomorphic-ope}; the
chain-level lift with the stated cotangent hypotheses is
Theorem~\ref{thm:app-unreduced-cotangent-lift-obstruction}.
\end{proof}

\begin{rmk}[Formal-Darboux stalk and global conjecture]
Theorem~\ref{app:thm:stalk-of-conj} computes the trace-sector
formal-Darboux fibre of Conjecture~\ref{app:conj:mixed-HT-Deligne} at the
formal-Darboux vacuum $b_p$.  The assertion that the global comparison on
$\R^2\times\C^2$ between the bulk
$\mathcal A^{\mathrm{cl}}_{\mathrm{bulk}}$ and the mixed-HT chiral
derived centre of the boundary open category at $d=2$ is a
quasi-isomorphism is
Conjecture~\ref{app:conj:mixed-HT-Deligne}.  The \(d=2\)
chiral \(E_d\)-Deligne obstruction is
\(\operatorname{Ob}^{\mathrm{ch}\,E_2}_{\mathrm{Deligne}}\)
(\S\ref{ssec:open-chiral-Ed-Deligne}).  Identification
with any curve chiral algebra requires a curve-restriction theorem and
matched-conventions transfer
(Theorem~\ref{thm:matched-conventions-transfer}).
\end{rmk}

\subsection{The chiral \texorpdfstring{$E_d$}{Ed}-operad: holomorphic little disks with principal-parts data}\label{app:hfc:chiral-Ed-operad}

The functorial Mixed-HT Deligne conjecture
(Conjecture~\ref{app:conj:mixed-HT-Deligne}) requires an operadic
action: the chiral derived centre
$Z^{\mathrm{der}}_{E_d^{\mathrm{ch}}}(\mathcal A)$ is the endomorphism
object in a chiral $E_d$-module category.  At $d=1$ the chiral operad
on curves gives this structure
\cite{beilinson-drinfeld-chiral}*{\S 3.4}.  At $d\geq 2$ the operad
is controlled by higher-codimension diagonal local cohomology.  In the
formal \(d=2\) Hamiltonian-current sector the all-arity principal-part
operations are the tree kernels of
Theorem~\ref{app:thm:formal-polydisk-all-arity-current-operad}.  The
remaining \(d=2\) chiral Deligne problem is not the existence of these
formal current kernels; it is the Maurer--Cartan lift in
\(\mathfrak{Def}^{\mathrm{HT}}_{\mathrm{Mor}\mbox{-}E_2}
(\mathcal C^{\mathrm{op}}_\partial,b_p)\), Ran descent, and
compatibility with the Costello QME complex.

\begin{defn}[Configuration $\mathcal D$-module]\label{app:defn:conf-D-module}
Let $X$ be a smooth complex variety of complex dimension $d$ and $I$
a non-empty finite set.  Write $X^I$ for the $|I|$-fold product and
$j_I\colon\mathrm{Conf}_I(X)\hookrightarrow X^I$ for the open embedding
of the configuration space of $|I|$ distinct points.  The
\emph{configuration $\mathcal D$-module}
\[
  \omega^{\mathrm{ch}}_{X^I}=j_{I,*}j_I^!\omega_{X^I}
\]
is the $!$-extension to $X^I$ of the dualising sheaf of
$\mathrm{Conf}_I(X)$.
\end{defn}

\begin{defn}[Chiral $E_d$ coefficient system]\label{app:defn:chiral-Ed-operad}
The \emph{chiral $E_d$ coefficient system} on $\C^d$ is the
collection \(\mathsf{hE}_d^{\mathrm{ch}}\) with single colour and
coefficient object at
arity $|I|$
\[
  \mathsf{hE}_d^{\mathrm{ch}}(I)
   =\omega^{\mathrm{ch}}_{(\C^d)^I},
\]
and with composition given by the chiral product on
\(!\)-pullback along the diagonal embedding
$\Delta_A\colon(\C^d)^A\hookrightarrow(\C^d)^I$ for a partition
$I=\bigsqcup_{a\in A}I_a$, following
\cite{beilinson-drinfeld-chiral}*{\S 3.4} as the formal
higher-codimension Ran-space analogue.  At \(d\geq2\) the action on
arbitrary mixed-HT factorization algebras, and its comparison with a
Ran \(\mathcal D\)-module model, is additional structure.  In the
formal \(d=2\) Hamiltonian-current model the tree-level all-arity
suboperad is constructed in
Theorem~\ref{app:thm:formal-polydisk-all-arity-current-operad}.
\end{defn}

\begin{rmk}[Comparison with curves at \texorpdfstring{$d=1$}{d=1}]\label{app:rmk:hEd-vs-BD-d1}
For $d=1$ on a curve $C$, the configuration $\mathcal D$-module
$\omega^{\mathrm{ch}}_{C^I}$ is the chiral operad on curves of
\cite{beilinson-drinfeld-chiral}*{\S 3.4}.  At $d\geq 2$ the diagonals
$\Delta_{ij}\subset(\C^d)^I$ have complex codimension $d$, so the
principal parts on $\omega^{\mathrm{ch}}_{(\C^d)^I}$ involve
distributions of $d$-fold polar order; this is the higher-codimension
chiral structure specific to \(d\geq2\).
\end{rmk}

\begin{rmk}[Arity two at \texorpdfstring{$d=2$}{d=2}]
For $X=\C^2$ and $I=\{1,2\}$, the singular quotient of
\(\omega^{\mathrm{ch}}_{X^I}\) along the pairwise diagonal is
\[
  H^2_\Delta(\C[[z_1,z_2,w_1,w_2]])\,d\zeta_1d\zeta_2
  =
  \bigoplus_{a,b\ge0}\C[[w_1,w_2]]\rho^\Delta_{a,b}.
\]
The generator \(\rho^\Delta_{0,0}=K^{(2)}_\Delta\) is the
two-dimensional Cauchy kernel.  For the Hamiltonian BF algebra, the
associated binary operation sends \(x\otimes y\) to
\([x,y]_{\mathfrak g}\), giving the formal singular product of
Theorem~\ref{thm:formal-two-dimensional-holomorphic-ope}.
\end{rmk}

\begin{defn}[Tree principal-part kernels on the formal surface]
\label{app:defn:tree-principal-part-kernels}
Let \(I\) be a finite set with \(|I|\geq2\).  A rooted binary tree
\(T\) with leaves labelled by \(I\) determines an iterated partial
diagonal
\[
  \Delta_T\subset(\widehat{\C^2}_0)^I
\]
by collapsing the two incoming clusters at each internal vertex.  For
each internal edge \(e\) let
\(\zeta_{e,1},\zeta_{e,2}\) be the two normal formal coordinates to the
corresponding collision divisor, chosen by subtracting the two cluster
coordinates.  The tree principal-part kernel is the local-cohomology
class
\[
  K_T
  =
  \bigwedge_{e\in E_{\mathrm{int}}(T)}
  \frac{d\zeta_{e,1}\,d\zeta_{e,2}}
       {\zeta_{e,1}\zeta_{e,2}}
  \in
  H^{2(|I|-1)}_{\Delta_T}
  \bigl(\mathcal O_{(\widehat{\C^2}_0)^I}\bigr)
  \otimes \omega_{(\C^2)^{I}/\C^2}.
\]
The relative volume form is taken modulo translation of the root
cluster.  Changing the cluster-coordinate splitting changes \(K_T\) by
a Cech coboundary in the local-cohomology complex, hence not in
cohomology.
\end{defn}

\begin{defn}[Formal \(d=2\) tree-kernel current operad]
\label{app:defn:tree-kernel-current-operad}
Let \(\mathsf{PP}^{(2)}_{\mathrm{tree}}\) be the dg operad whose
arity-\(I\) component is spanned by the classes \(K_T\) of
Definition~\ref{app:defn:tree-principal-part-kernels}, for rooted
binary \(I\)-trees \(T\), modulo:
\begin{enumerate}[label=\textup{(\roman*)}]
  \item antisymmetry under interchange of the two incoming branches at
        a binary vertex;
  \item the Arnold--Jacobi relation on every three-leaf subtree,
        \[
          \eta_{12}\wedge\eta_{23}
          +\eta_{23}\wedge\eta_{31}
          +\eta_{31}\wedge\eta_{12}=0,
          \qquad
          \eta_{ij}=
          \frac{d(z_{i,1}-z_{j,1})\,d(z_{i,2}-z_{j,2})}
          {(z_{i,1}-z_{j,1})(z_{i,2}-z_{j,2})};
        \]
  \item Cech coboundaries in the local-cohomology representatives.
\end{enumerate}
Partial composition grafts rooted trees, pulls back the two kernel
classes to the iterated diagonal, and wedges the forms.  The unit is
the empty collision kernel in arity one.
\end{defn}

\begin{lem}[Arity-three residue and Jacobi test]
\label{app:lem:arity-three-residue-jacobi}
  Let \(z^{(1)},z^{(2)},z^{(3)}\) be three formal points of
  \(\widehat{\C^2}_0\), and put
  \[
    a=z^{(1)}-z^{(2)},\qquad b=z^{(2)}-z^{(3)}.
  \]
  The small diagonal is the complete intersection \(a=b=0\).  For every
  formal Laurent principal part \(F(a,b,z^{(3)})\),
  \[
    \operatorname{Res}_{a=0}\operatorname{Res}_{b=0}
    \left(
      F\,
      \frac{da_1\,da_2}{a_1a_2}
      \wedge
      \frac{db_1\,db_2}{b_1b_2}
    \right)
    =
    \operatorname{Res}_{b=0}\operatorname{Res}_{a=0}
    \left(
      F\,
      \frac{da_1\,da_2}{a_1a_2}
      \wedge
      \frac{db_1\,db_2}{b_1b_2}
    \right).
  \]
  Both sides are the Grothendieck residue of \(F\) on
  \(\Delta_{\{1,2,3\}}\).  Hence the two nested-residue orders for the
  three-leaf tree kernel define the same local-cohomology class.

  For homogeneous \(x_1,x_2,x_3\in\mathfrak g\), the arity-three
  operation sends the signed Jacobi generator
  \[
    \sum_{\mathrm{cyc}(1,2,3)}
    (-1)^{|x_i||x_k|}
    \bigl(\eta_{ij}\wedge\eta_{jk}\bigr)
    \otimes [[x_i,x_j],x_k]
  \]
  to zero in the quotient action of
  \(\mathsf{PP}^{(2)}_{\mathrm{tree}}\).  The kernel component is the
  Arnold local-cohomology relation, and the coefficient component is the
  graded Jacobi identity in the degree-zero dg Lie bracket of
  \(\mathfrak g\).  The cyclic sum runs over
  \((i,j,k)=(1,2,3),(2,3,1),(3,1,2)\).
\end{lem}

\begin{proof}
  The formal local ring at the small diagonal is
  \(\C[[a_1,a_2,b_1,b_2,z^{(3)}_1,z^{(3)}_2]]\).  The residue above is
  the coefficient functional extracting the coefficient of
  \(a_1^{-1}a_2^{-1}b_1^{-1}b_2^{-1}\).  Coefficient extraction in the
  \(a\)-variables commutes with coefficient extraction in the
  \(b\)-variables.  This proves residue Fubini and identifies both
  iterated residues with the Grothendieck residue for the regular
  sequence \(a_1,a_2,b_1,b_2\).

  The second assertion is the first operadic relation.  The three
  products \(\eta_{ij}\wedge\eta_{jk}\) are the three Cech
  representatives for the codimension-four collision obtained from the
  three cyclic two-step collisions.  Their signed sum is zero in the
  local-cohomology quotient by the Arnold relation of
  Definition~\ref{app:defn:tree-kernel-current-operad}.  The remaining
  coefficient is the signed Jacobiator of a degree-zero graded Lie
  bracket.  Therefore the arity-three relation is killed by the action.
\end{proof}

\begin{thm}[Formal-polydisk all-arity Hamiltonian current operad]
\label{app:thm:formal-polydisk-all-arity-current-operad}
Let
\[
  \mathfrak g=\mathfrak h\ltimes\mathfrak h^\vee_{\mathrm{cont}}[1],
  \qquad
  \mathfrak h=\C[[z_1,z_2]]/\C\cdot1,
\]
with the residue pairing and continuous bracket of the formal
Hamiltonian BF theory.  The operad
\(\mathsf{PP}^{(2)}_{\mathrm{tree}}\) acts on the formal Hamiltonian
current factorization algebra by
\[
  \mu_T(K_T;x_1,\ldots,x_{|I|})
  =
  K_T\,\mathsf{J}_{[-,\ldots,-]_T(x_1,\ldots,x_{|I|})},
  \qquad x_i\in\mathfrak g,
\]
where \([-,\ldots,-]_T\) is the iterated bracket prescribed by the
rooted tree \(T\).  In arity two this is the diagonal singular product
\[
  \mathsf{J}_x(z)\mathsf{J}_y(w)
  \sim
  \frac{d\zeta_1\,d\zeta_2}{\zeta_1\zeta_2}
  \mathsf{J}_{[x,y]_{\mathfrak g}}(w)
\]
of Theorem~\ref{thm:formal-two-dimensional-holomorphic-ope}.  The
partial compositions are associative on cohomology, and the action is
compatible with the Chevalley--Eilenberg differential.
\end{thm}

\begin{proof}
The local-cohomology class of an iterated collision is computed by
successive Grothendieck residues along the normal coordinates
\(\zeta_{e,1},\zeta_{e,2}\).  Lemma~\ref{app:lem:arity-three-residue-jacobi}
is the first nontrivial case: residue Fubini identifies the two nested
orders on the small diagonal, and the three-leaf Jacobi generator is
killed by the Arnold local-cohomology relation together with the graded
Jacobi identity in \(\mathfrak g\).  For a larger tree, every
associativity comparison is a sequence of such three-leaf rotations
and independent residue interchanges.  Antisymmetry of the kernel and
antisymmetry of the
Lie bracket agree at each binary vertex.  The Chevalley--Eilenberg
differential is the derivation dual to the same bracket and is
continuous in the Tate topology; hence it commutes with the operation
maps.  The arity-two formula is the definition of the generator
\(K^{(2)}_\Delta\), and every \(K_T\) is obtained from it by operadic
composition.
\end{proof}

\begin{rmk}[What this theorem does not supply]
\label{app:rmk:tree-kernel-not-global-deligne}
Theorem~\ref{app:thm:formal-polydisk-all-arity-current-operad}
constructs the all-arity formal current-kernel layer.  It does not by
itself construct the Morita chiral derived centre of the open boundary
category, a Ran \(\mathcal D\)-module presentation,
Costello--Gwilliam descent on the brane-link boundary, or a Costello
QME solution.  Those are the remaining components of
\(\operatorname{Ob}^{\mathrm{ch}\,E_2}_{\mathrm{Deligne}}\),
\(\operatorname{Ob}^{\mathrm{desc}}_{\partial}\), and
\(\operatorname{Ob}^{\mathrm{QME}}_{\mathrm{all\mbox{-}order}}\).
\end{rmk}

\begin{constr}[$\mathsf{hE}_d^{\mathrm{ch}}$ operation maps under the operadic action]\label{app:prop:hEd-algebras}
Assume the chiral \(E_d\)-Deligne operadic datum of
\S\ref{ssec:open-chiral-Ed-Deligne} is supplied on
\(\mathrm{HolFA}_d(\C^d)\).  Then a holomorphic factorization
\(E_d\)-algebra
\(\mathcal A\in\mathrm{HolFA}_d(\C^d)\)
(Definition~\ref{app:defn:hol-fact-Ed}) carries structure maps
\[
  \mu^{\mathrm{ch}}_I\colon
   \mathsf{hE}_d^{\mathrm{ch}}(I)\otimes^!\mathcal A_{\mathrm{Ran}}^{\boxtimes I}
   \longrightarrow\Delta_I^*\mathcal A_{\mathrm{Ran}}
\]
intertwining the chiral product of
Definition~\ref{app:defn:chiral-Ed-operad} with $\Delta_I^*$, where
\(\mathcal A_{\mathrm{Ran}}\) denotes the corresponding Ran-space
factorization object.  At \(d=1\) these maps are the chiral operations
on curves.  At \(d=2\), the formal-Darboux Hamiltonian-current
coefficients are the all-arity tree kernels of
Theorem~\ref{app:thm:formal-polydisk-all-arity-current-operad}.
\end{constr}

\begin{rmk}[Operation maps]
The Costello--Gwilliam and Lurie--Ayala--Francis factorization
formalisms give the symmetric monoidal factorization category and its
factorization-homology presentation.  The chiral \(E_d\)-Deligne
equivalence at \(d\geq2\) is the theorem that these operation maps lift
to the Morita chiral centre through the relative Maurer--Cartan complex
and descent.  Construction~\ref{app:prop:hEd-algebras}
is the global model; the proved formal-current part at \(d=2\) is
Theorem~\ref{app:thm:formal-polydisk-all-arity-current-operad}, and the
relative lift is controlled by
Definition~\ref{app:defn:morita-chiral-E2-extension-complex}.
\end{rmk}

\begin{constr}[Mixed-HT enhancement under the operadic action
\texorpdfstring{$\mathsf{hE}_d^{\mathrm{ch},\mathrm{HT}}$}{hEd-ch-HT}]\label{app:constr:hEd-HT}
Adjoin to the system $\mathsf{hE}_d^{\mathrm{ch}}$ the
topological \(E_k\)-coefficients on $\R^k_{\mathrm{top}}$.  The
resulting mixed-HT coefficient system
\(\mathsf{hE}_d^{\mathrm{ch},\mathrm{HT}}\) has arity-\(I\)
coefficient
\[
  \mathsf{hE}_d^{\mathrm{ch},\mathrm{HT}}(I)
   =\mathrm{Conf}_I(\R^k)\times
   \omega^{\mathrm{ch}}_{(\C^d)^I},
\]
with composition the product of topological little-disks
composition on the $\R^k$ factor and the chiral composition of
Definition~\ref{app:defn:chiral-Ed-operad} on the $\C^d$ factor.  An
action of this system on a mixed-HT factorization algebra is an
all-arity operadic structure.
\end{constr}

\begin{defn}[Morita chiral \(E_2\)-extension complex]
\label{app:defn:morita-chiral-E2-extension-complex}
Let
\[
  \mathsf Q=\mathsf{PP}^{(2)}_{\mathrm{tree}},
  \qquad
  \mathsf P=\mathsf{hE}_2^{\mathrm{ch},\mathrm{HT}}
\]
with \(\mathsf Q\) restricted to the formal Hamiltonian-current sector on
\(\widehat{\C^2}_0\).  Choose cofibrant dg-operad models for
\(\mathsf P\) and \(\mathsf Q\), and write
\(\mathrm B\mathsf P\), \(\mathrm B\mathsf Q\) for their bar
cooperads.  Put
\[
  Z_\partial
  =
  Z^{\mathrm{der},\mathrm{Mor}}_{E_2^{\mathrm{HT}}}
    (\mathcal C^{\mathrm{op}}_\partial).
\]
Evaluation at the formal-Darboux brane \(b_p\), followed by scalar
reduction to the trace-sector chart of
Theorem~\ref{app:thm:stalk-of-conj}, gives a restriction morphism of
complete filtered convolution dg Lie algebras
\[
  \operatorname{res}_{b_p}\colon
  \mathrm{Conv}\bigl(\mathrm B\mathsf P,\operatorname{End}_{Z_\partial}\bigr)
  \longrightarrow
  \mathrm{Conv}\bigl(\mathrm B\mathsf Q,
      \operatorname{End}_{\operatorname{ev}_{b_p}(Z_\partial)^{\mathrm{tr}}}
      \bigr).
\]
The \emph{Morita chiral \(E_2\)-extension complex} is the shifted
homotopy fibre
\[
  \mathfrak{Def}^{\mathrm{HT}}_{\mathrm{Mor}\mbox{-}E_2}
    (\mathcal C^{\mathrm{op}}_\partial,b_p)
  :=
  \operatorname{fib}(\operatorname{res}_{b_p})[-1],
\]
with differential twisted by the Maurer--Cartan element corresponding to
the formal tree-kernel current action of
Theorem~\ref{app:thm:formal-polydisk-all-arity-current-operad}.
The class
\[
  \operatorname{Ob}^{\mathrm{ch}\,E_2}_{\mathrm{Deligne}}
  \in
  H^2\!\left(
    \mathfrak{Def}^{\mathrm{HT}}_{\mathrm{Mor}\mbox{-}E_2}
    (\mathcal C^{\mathrm{op}}_\partial,b_p)
  \right)
\]
is the first obstruction to lifting the formal current-kernel
\(\mathsf Q\)-action to a Morita mixed-HT chiral \(E_2\)-centre action on
\(Z_\partial\).
\end{defn}

\begin{prop}[Maurer--Cartan control of the Morita centre lift]
\label{app:prop:morita-centre-lift-MC-control}
Assume the filtration on
\(\mathfrak{Def}^{\mathrm{HT}}_{\mathrm{Mor}\mbox{-}E_2}
(\mathcal C^{\mathrm{op}}_\partial,b_p)\) is complete and pronilpotent.
A Morita mixed-HT chiral \(E_2\)-centre action on
\(Z_\partial\) extending the formal tree-kernel current action is
equivalent to a Maurer--Cartan lift along
\(\operatorname{res}_{b_p}\).  If such a lift is constructed modulo
filtration \(F^{r+1}\), the obstruction to extending it modulo
\(F^{r+2}\) is a class in
\[
  H^2\!\left(
    \operatorname{gr}^{r+1}
    \mathfrak{Def}^{\mathrm{HT}}_{\mathrm{Mor}\mbox{-}E_2}
      (\mathcal C^{\mathrm{op}}_\partial,b_p)
  \right).
\]
The vanishing of all these classes, with chosen degree-one primitives,
is exactly the operadic part of the global mixed-HT Deligne datum.
\end{prop}

\begin{proof}
An algebra over a cofibrant dg operad \(\mathsf R\) on a complex \(V\) is
a Maurer--Cartan element of
\(\mathrm{Conv}(\mathrm B\mathsf R,\operatorname{End}_V)\).  Restriction
from \(\mathsf P\) to the formal tree-kernel sector \(\mathsf Q\), and
then to the \(b_p\)-Hochschild chart, is the morphism
\(\operatorname{res}_{b_p}\).  Therefore the extension problem is the
standard lifting problem for Maurer--Cartan elements in the homotopy
fibre of this morphism.  Completeness and pronilpotence give the usual
filtration-by-filtration obstruction theory: the quadratic term of the
Maurer--Cartan equation at order \(r+1\) is closed in the associated
graded complex, changing the partial lift changes it by a coboundary,
and a degree-one primitive is precisely the next correction term.
\end{proof}

\subsection{Three hypotheses for chiral
\texorpdfstring{$E_d$}{Ed} Deligne at \texorpdfstring{$d\geq 2$}{d>=2}}\label{app:hfc:three-hypotheses}

The chiral \(E_d\) Deligne problem at \(d\geq2\) decomposes here into
three obstruction classes for the coefficient system
\(\mathsf{hE}_d^{\mathrm{ch}}\) of
Definition~\ref{app:defn:chiral-Ed-operad}.  In the formal
Hamiltonian-current sector at \(d=2\) the all-arity tree-kernel
suboperad is constructed.  The remaining classes are the promotion to
the full chiral coefficient system, the Morita centre
Maurer--Cartan lift, and Ran descent.

\begin{defn}[Hypothesis F: formality of $\mathsf{hE}_d^{\mathrm{ch}}$]\label{app:hyp:formality}
The system \(\mathsf{hE}_d^{\mathrm{ch}}\), beyond the formal
tree-kernel suboperad of
Theorem~\ref{app:thm:formal-polydisk-all-arity-current-operad}, is
promoted to an all-arity \(\infty\)-operad, and that operad is formal:
there is
a quasi-isomorphism
\[
  \mathsf{hE}_d^{\mathrm{ch}}\;\simeq\;H_*(\mathsf{hE}_d^{\mathrm{ch}})
\]
to its homology operad in chain complexes of
$\mathcal D_{\mathrm{Ran}(\C^d)}$-modules.  At $d=1$ this is
Tamarkin's formality theorem
\cite{tamarkin-deligne}\cite{kontsevich-soibelman-deformation}; at
$d\geq 2$ the obstruction is the deformation class
\(\mathrm{Ob}_F=[\mathfrak m_d]\).
\end{defn}

\begin{defn}[Hypothesis C: chiral centrality]\label{app:hyp:centrality}
For $\mathcal A\in\mathrm{HolFA}_d(\C^d)$, there is a chiral bimodule
\(\infty\)-category
\(\mathcal A\otimes\mathcal A^{\mathrm{op}}\text{-Mod}^{\mathrm{ch}}\)
and an endomorphism object
\[
  Z^{\mathrm{der}}_{\mathsf{hE}_d^{\mathrm{ch}}}(\mathcal A)
   =\RHom_{\mathcal A\otimes\mathcal A^{\mathrm{op}}}(\mathcal A,\mathcal A).
\]
This object acts on \(\mathcal A\) by homotopy-central chiral
operations compatible with factorization products.
At \(d=1\) this is the chiral-centre construction for chiral algebras
on curves.
For \(d\geq2\) its topological analogue is Lurie's
\(E_d\)-centre construction \cite{lurie-ha}*{\S 5.3.1}; the
holomorphic chiral enhancement is governed by the corresponding
operadic obstruction problem.
\end{defn}

\begin{defn}[Hypothesis R: Ran \( ! \)-descent]\label{app:hyp:ran-descent}
There is a chosen Ran \(\mathcal D\)-module presentation
\(\mathcal A_{\mathrm{Ran}}\) of the Costello--Gwilliam
factorization algebra \(\mathcal A\), and \(\mathcal A_{\mathrm{Ran}}\)
satisfies \( ! \)-descent on the Ran site.  Equivalently, the
comparison from the Costello--Gwilliam Weiss descent model to the Ran
\( ! \)-descent model is a quasi-isomorphism in the selected
coefficient category.  At \(d=1\) this is the Beilinson--Drinfeld
factorization theorem for chiral algebras on curves
\cite{beilinson-drinfeld-chiral}*{\S 3}.  At \(d\geq2\) it is a
hypothesis of the global chiral \(E_d\)-Deligne theorem.
\end{defn}

\begin{thm}[Chiral $E_d$ Deligne under formality and Ran descent]\label{app:thm:relative-Ed-deligne}
Under Hypotheses~\ref{app:hyp:formality}, \ref{app:hyp:centrality},
\ref{app:hyp:ran-descent}, the chiral derived centre satisfies
\[
  Z^{\mathrm{der}}_{\mathsf{hE}_d^{\mathrm{ch}}}(\mathcal A)
   \;\simeq\;
   C^\bullet_{\mathrm{ch}}(\mathcal A,\mathcal A)
\]
as $E_{d+1}^{\mathrm{ch}}$-algebras.  Hypotheses~\ref{app:hyp:centrality}
and~\ref{app:hyp:ran-descent} are the chiral enhancements of known
topological and curve-level theorems; at \(d\geq2\) they remain
hypotheses of the chiral \(E_d\)-Deligne statement together with
Hypothesis~\ref{app:hyp:formality}.
\end{thm}

\begin{proof}
By Hypothesis~\ref{app:hyp:centrality} the chiral bimodule
\(\infty\)-category
\(\mathcal A\otimes\mathcal A^{\mathrm{op}}\text{-Mod}^{\mathrm{ch}}\)
exists and the derived centre is the endomorphism object of the
diagonal bimodule,
\(Z^{\mathrm{der}}_{\mathsf{hE}_d^{\mathrm{ch}}}(\mathcal A)
=\RHom_{\mathcal A\otimes\mathcal A^{\mathrm{op}}}(\mathcal A,\mathcal A)\).

\emph{The chiral two-sided bar resolution.}  Write
\(\mu_{\mathrm{ch}}\colon j_*j^*(\mathcal A\boxtimes\mathcal A)\to
\Delta_!\mathcal A\) for the chiral product.  The two-sided bar complex
\[
  \mathrm{Bar}_\bullet(\mathcal A)
  =\Bigl(\mathcal A\otimes^!\mathcal A^{\otimes^!\bullet}
    \otimes^!\mathcal A,\;b\Bigr),
  \qquad
  b=\sum_i(-1)^i\,\mu_{\mathrm{ch}}^{(i,i+1)},
\]
with \(b\) the alternating sum of chiral multiplications of adjacent
factors, resolves the diagonal bimodule \(\mathcal A\) by \(j\)-free
\(\mathcal A\)-bimodules on \(\operatorname{Ran}(X)\): the extra
degeneracy \(s_{-1}=\eta\otimes\mathrm{id}\) built from the unit
\(\eta\) gives a chiral contraction \(bs_{-1}+s_{-1}b=\mathrm{id}\),
splitting the augmentation, so
\(\mathrm{Bar}_\bullet(\mathcal A)\to\mathcal A\) is a
quasi-isomorphism, and each \(\mathrm{Bar}_n\) is induced from the free
bimodule on \(\mathcal A^{\otimes^! n}\).  Hence
\[
  \RHom_{\mathcal A\otimes\mathcal A^{\mathrm{op}}}(\mathcal A,\mathcal A)
  =\operatorname{Hom}_{\mathcal A\text{-bimod}}\!\bigl(
   \mathrm{Bar}_\bullet(\mathcal A),\mathcal A\bigr)
  =\operatorname{Tot}\operatorname{Hom}\!\bigl(
   \mathcal A^{\otimes^!\bullet},\mathcal A\bigr)
  =:C^\bullet_{\mathrm{ch}}(\mathcal A,\mathcal A),
\]
the chiral Hochschild cochain complex with the differential dual to
\(b\); at \(d=1\) this is the Beilinson--Drinfeld chiral Hochschild
complex \cite{beilinson-drinfeld-chiral}*{\S3}.

\emph{Model independence.}  By Hypothesis~\ref{app:hyp:ran-descent} the
Ran \(\mathcal D\)-module presentation satisfies \(!\)-descent and its
comparison to the Costello--Gwilliam Weiss model is a
quasi-isomorphism, so \(\mathrm{Bar}_\bullet(\mathcal A)\) and the
resulting \(C^\bullet_{\mathrm{ch}}(\mathcal A,\mathcal A)\) are the same
object of \(\mathcal B^{\mathrm{adm}}\) in either model.

\emph{The \(E_{d+1}^{\mathrm{ch}}\)-structure.}  The Hochschild complex
carries the chiral brace operations: a tuple of cochains is composed by
inserting outputs into inputs along the chiral product, the chiral
analogue of the McClure--Smith brace action of the little-\(d\)-cubes
chains, assembling \(C^\bullet_{\mathrm{ch}}\) into an
\(\mathsf{hE}_d^{\mathrm{ch}}\)-algebra.  By
Hypothesis~\ref{app:hyp:formality} the operad quasi-isomorphism
\(\mathsf{hE}_d^{\mathrm{ch}}\simeq
H_*(\mathsf{hE}_d^{\mathrm{ch}})=\mathsf P_{d+1}^{\mathrm{ch}}\) to the
chiral Poisson operad transfers the brace action to a
\(\mathsf P_{d+1}^{\mathrm{ch}}\)-, hence
\(E_{d+1}^{\mathrm{ch}}\)-, algebra structure on the same complex --
the chiral Deligne structure.

\emph{Conclusion.}  The underlying complexes agree by the bar resolution
and model independence, and both carry the transferred
\(E_{d+1}^{\mathrm{ch}}\)-structure compatibly with the chiral product;
hence \(Z^{\mathrm{der}}_{\mathsf{hE}_d^{\mathrm{ch}}}(\mathcal A)\simeq
C^\bullet_{\mathrm{ch}}(\mathcal A,\mathcal A)\) as
\(E_{d+1}^{\mathrm{ch}}\)-algebras.
\end{proof}

For the formal-Darboux Hamiltonian BF model the residue-kernel operad
gives the formal current-kernel input.  It is not, by itself, a proof of
formality for the full chiral \(\mathsf hE_2^{\mathrm{ch}}\)-operad, a
Ran descent theorem, or a Morita centre lift.  The next theorem records
the exact criterion that promotes this formal current layer to the
stalkwise chiral \(E_2\)-Deligne statement.

\begin{thm}[Formal current-kernel \(E_2\) criterion at the Hamiltonian BF stalk]\label{app:thm:FCR-at-stalk}
Let \(\mathcal A=\mathcal A^{\mathrm{cl}/\hbar}_{\mathrm{bulk}}\) be the
formal-Darboux Hamiltonian BF factorization algebra of
\(\mathfrak g=\mathfrak h\ltimes\mathfrak h^\vee_{\mathrm{cont}}[1]\) at
a Dirac brane stack, with the brane-preserving \(\Omega\)-background
grading of Definition~\ref{defn:omega-grading} and nonresonant
\((\epsilon_1,\epsilon_2)\).  The tree-kernel operad
\(\mathsf{PP}^{(2)}_{\mathrm{tree}}\) acts on the formal Hamiltonian
current sector by Theorem~\ref{app:thm:formal-polydisk-all-arity-current-operad}.
The assertion that this action extends to a stalkwise chiral
\(E_2\)-Deligne equivalence is equivalent to the following additional
data:
\[
  \mathrm{Ob}^{\mathrm{stalk}}_{\mathrm{ch}E_2}
  =
  \left(
    \mathrm{Ob}_{F}^{\mathrm{stalk}},
    \operatorname{Ob}^{\mathrm{ch}\,E_2}_{\mathrm{Deligne}},
    \mathrm{Ob}_{R}^{\mathrm{stalk}},
    \{\operatorname{res}^{0}_{\mathrm{nsc}}(\mathfrak o_n)\}_{n\ge2}
  \right)=0
\]
with chosen primitives.  Here
\(\mathrm{Ob}_{F}^{\mathrm{stalk}}\) is the obstruction to promoting
\(\mathsf{PP}^{(2)}_{\mathrm{tree}}\) to the chosen cofibrant
\(\mathsf hE_2^{\mathrm{ch},\mathrm{HT}}\)-model,
\(\operatorname{Ob}^{\mathrm{ch}\,E_2}_{\mathrm{Deligne}}\) is the
Morita centre-lift class of
Definition~\ref{app:defn:morita-chiral-E2-extension-complex},
\(\mathrm{Ob}_{R}^{\mathrm{stalk}}\) is the Ran/Weiss comparison class,
and
\[
  \operatorname{res}^{0}_{\mathrm{nsc}}(\mathfrak o_n)
  \in H^2(\ker\pi_{\mathrm{sc}},Q)
\]
is the residual non-scalar weight-zero centrality/QME class of
Lemma~\ref{lem:nonscalar-weight-zero-vanishing}.  On this vanishing
locus, with the displayed primitives chosen,
Theorem~\ref{app:thm:relative-Ed-deligne} gives
\[
  Z^{\mathrm{der}}_{\mathsf{hE}_2^{\mathrm{ch}}}(\mathcal A)
   \;\simeq\;
   C^\bullet_{\mathrm{ch}}(\mathcal A,\mathcal A)
\]
as \(E_3^{\mathrm{ch}}\)-algebras, projectively with scalar Capelli
curvature \(\hbar N[\bar c]\).
\end{thm}

\begin{proof}
\emph{Formal current kernels.}  The all-arity current action of
Theorem~\ref{app:thm:formal-polydisk-all-arity-current-operad} is by the
residue tree kernels \(K_T\) of
Definition~\ref{app:defn:tree-principal-part-kernels}, which compose
strictly associatively by Grothendieck residue Fubini
(Lemma~\ref{app:lem:arity-three-residue-jacobi}) and are \(Q\)-closed by
the Chevalley--Eilenberg identity.  Every higher singular product is an
iterated binary residue collision
(Corollary~\ref{cor:binary-generation-higher-two-dimensional-ope}), so
the formal Hamiltonian-current suboperad is strict.

\emph{Promotion data.}  A lift from the formal tree-kernel suboperad to
the chosen \(\mathsf hE_2^{\mathrm{ch},\mathrm{HT}}\)-model is a
Maurer--Cartan lifting problem; its first obstruction is
\(\mathrm{Ob}_{F}^{\mathrm{stalk}}\).  A homotopy-central Morita action
on \(Z_\partial\) is controlled by
Proposition~\ref{app:prop:morita-centre-lift-MC-control}, with
obstruction class
\(\operatorname{Ob}^{\mathrm{ch}\,E_2}_{\mathrm{Deligne}}\).  The
Ran/Weiss comparison is the class
\(\mathrm{Ob}_{R}^{\mathrm{stalk}}\).  The
\(\Omega\)-localization theorem contracts only nonzero weights and
computes the scalar weight-zero projection as the Capelli class; the
non-scalar weight-zero part is removable exactly when the residual
classes of Lemma~\ref{lem:nonscalar-weight-zero-vanishing} vanish and
compatible primitives are chosen.

If the four displayed obstruction components vanish, Hypotheses
\ref{app:hyp:formality}, \ref{app:hyp:centrality}, and
\ref{app:hyp:ran-descent} hold for this stalk with specified data.
Theorem~\ref{app:thm:relative-Ed-deligne} then gives the
\(E_3^{\mathrm{ch}}\)-equivalence, and its scalar weight-zero curvature
is the Capelli class by Theorem~\ref{thm:omega-localization-collapse}.
Conversely, such an \(E_3^{\mathrm{ch}}\)-equivalence restricts to the
formal current action, supplies the Maurer--Cartan lift, identifies the
Ran/Weiss model, and gives the residual centrality primitives; hence all
four obstruction components vanish.
The general \(d\ge2\) statement for an arbitrary holomorphic
factorization algebra remains Conjecture~\ref{conj:chiral-Ed-deligne}.
\end{proof}

The formal-Darboux stalk does not globalise by itself.  A Stein target
removes only the sheaf-cohomological part of a specified
\v Cech--Weiss descent problem; it does not construct the Morita centre
lift, the quantum master equation counterterms, the anomaly and locality
null-homotopies, or the comparison-cone contraction.

\begin{thm}[Stein descent criterion for a projective mixed-HT open--closed comparison]\label{app:thm:stein-global-deligne}
Let \(Y\) be a Stein holomorphic-symplectic surface,
\(X=\R^2_{\mathrm{top}}\times Y\), and \(L\subset X\) a Dirac brane
defect.  Choose a collared formal-Darboux Weiss cover
\(\mathfrak U=\{U_a\}_{a\in A}\) of the brane-link boundary.  Suppose the
following data are fixed.
\begin{enumerate}[label=\textup{(S\arabic*)}]
\item On every \(U_a\), the formal current-kernel criterion of
  Theorem~\ref{app:thm:FCR-at-stalk} holds: the formality,
  Morita-centre, Ran/Weiss, and residual weight-zero obstruction classes
  vanish and compatible primitives are chosen.
\item The local boundary categories form a Weiss cosheaf, and the local
  mapping complexes
  \[
    \operatorname{Hom}_{\mathcal B^{\mathrm{adm}}}
    \bigl(\mathcal A^{\mathrm{cl}/\hbar}_{\mathrm{bulk}}|_{U_a},
    Z^{\mathrm{der},\mathrm{Mor}}_{E_2^{\mathrm{HT}}}
      (\mathcal C^{\mathrm{op}}_\partial)|_{U_a}\bigr)
  \]
  are represented by nuclear Fr\'echet Dolbeault complexes satisfying
  Cartan--H\"ormander acyclicity on the finite intersections of
  \(\mathfrak U\).
\item The local morphisms \(F_a\), overlap homotopies \(H_{ab}\), and
  higher cochains give a homotopy-coherent null-homotopy of
  \[
    \operatorname{Ob}_{\mathrm{desc}}\in
    H^2\!\left(
    \operatorname{Tot}\check C^\bullet_{\mathrm{Weiss}}
    (\mathfrak U;\mathcal E^\bullet_{\mathrm{desc}})
    \right)
  \]
  and of its higher obstruction tower
  (Definition~\ref{app:defn:descent-obstruction-complex}).
\item The all-scale quantum master equation datum, anomaly and locality
  null-homotopies, and comparison-cone contraction of
  Definition~\ref{defn:oca-comparison-cone} are supplied.
\end{enumerate}
Then the local maps glue to a projective open--closed quasi-isomorphism
\[
  \mathrm{OCA}_N\colon
  \mathcal A^{\mathrm{cl}/\hbar}_{\mathrm{bulk}}
  \xrightarrow{\ \sim\ }
  Z^{\mathrm{der},\mathrm{Mor}}_{E_2^{\mathrm{HT}}}
    (\mathcal C^{\mathrm{op}}_\partial)
\]
in the projective category for the Capelli central extension.  Its
trace-sector scalar curvature is the projective Lie anomaly
\(\hbar N[\bar c]\).  The criterion applies to \(\C^2\), to the
analytifications of smooth affine holomorphic-symplectic surfaces, and
to cotangent bundles \(T^*\Sigma\) of affine curves when
\textup{(S1)--(S4)} are verified.  The punctured plane
\(\C^2\setminus\{0\}\) is not covered by the Stein conclusion.
\end{thm}

\begin{proof}
\emph{Local data.}  Condition \textup{(S1)} is exactly the obstruction
criterion of Theorem~\ref{app:thm:FCR-at-stalk}.  On each member of the
Weiss cover it gives a projective local morphism \(F_a\) with scalar
curvature \(\hbar N[\bar c]\); it also records the chosen primitives
needed for the Morita centre lift and residual weight-zero terms.  This
is local formal-Darboux data, not a global morphism.

\emph{\v Cech--Weiss descent.}  Conditions \textup{(S2)} and
\textup{(S3)} place the local \(F_a\)'s in the total complex
\(\operatorname{Tot}\check C^\bullet_{\mathrm{Weiss}}
(\mathfrak U;\mathcal E^\bullet_{\mathrm{desc}})\) of
Definition~\ref{app:defn:descent-obstruction-complex}.  Cartan--H\"ormander
acyclicity on Stein finite intersections kills only the ambient positive
\v Cech cohomology of the specified Dolbeault mapping complexes.  The
actual descent obstruction is the class of \(\delta_{\check C}H\) and
its higher coherences; by \textup{(S3)} those classes have specified
null-homotopies.  Therefore the \(F_a\)'s assemble to a global
homotopy-coherent open--closed morphism.

\emph{Quantum and cone data.}  Condition \textup{(S4)} supplies the QME
counterterms, the anomaly/locality null-homotopies, and the contraction
of the global comparison cone.  Hence the assembled morphism is a
quasi-isomorphism in the projective category.  The scalar part of its
curvature is unchanged by the descent homotopies, because on every
chart it is the central Capelli class \(\hbar N[\bar c]\) of
Theorem~\ref{thm:capelli-omega-central-identification}.
\end{proof}

\begin{rmk}[Compact targets and the period twist]\label{app:rmk:compact-period-twist}
On a compact holomorphic-symplectic surface the Darboux cover is no
longer Leray for \(\C_X\), and the Hamiltonian period
\(\operatorname{per}_X(v)=[h_i-h_j]\in H^1(X,\C_X)\) of
\S\ref{ssec:master-deformation-local-datum} becomes an additional
descent obstruction.  If \(\operatorname{per}_X(v)=0\), the analogue of
Theorem~\ref{app:thm:stein-global-deligne} still requires the same
local centre-lift, QME, anomaly/locality, and comparison-cone data.  A
nonzero period twists \(\mathrm{OCA}_N\) by the corresponding
\(\C^*\)-gerbe, and the remaining content is the chiral
\(E_2\)-Deligne comparison of Conjecture~\ref{conj:chiral-Ed-deligne}.
\end{rmk}

\begin{rmk}[Obstruction class for Hypothesis~\ref{app:hyp:formality}]\label{app:rmk:formality-obstruction}
Hypothesis~\ref{app:hyp:formality} is the assertion that the
Maurer--Cartan element
$\mathfrak m_d\in\mathrm{Def}\bigl(\mathsf{hE}_d^{\mathrm{ch}}\to
H_*(\mathsf{hE}_d^{\mathrm{ch}})\bigr)$ classifying the chiral
\(E_d\)-operad relative to its homology operad is gauge-equivalent to
zero.  At \(d=1\) Tamarkin~\cite{tamarkin-deligne} proves the
corresponding formality statement by associator methods.  At \(d\geq2\)
the corresponding chiral formality theorem is an independent datum.  The class
\[
  \mathrm{Ob}_F=[\mathfrak m_d]
   \in H^*\bigl(\mathrm{Def}(\mathsf{hE}_d^{\mathrm{ch}})\bigr)
\]
is the formality obstruction at $d\geq 2$.
\end{rmk}

\begin{rmk}[Chiral \(E_d\) obstruction]\label{app:rmk:chiral-Ed-obstruction}
The global chiral \(E_d\)-Deligne hypothesis at \(d\geq2\) consists of the
operadic action of Definition~\ref{app:defn:chiral-Ed-operad} together
with Hypotheses~\ref{app:hyp:formality}, \ref{app:hyp:centrality}, and
\ref{app:hyp:ran-descent}.  At the \(d=2\) Dirac-brane
formal-Darboux vacuum the Morita centre-action component is represented
by the extension complex
\(\mathfrak{Def}^{\mathrm{HT}}_{\mathrm{Mor}\mbox{-}E_2}
(\mathcal C^{\mathrm{op}}_\partial,b_p)\) of
Definition~\ref{app:defn:morita-chiral-E2-extension-complex}.  The full
obstruction is the class
\[
  \mathrm{Ob}_{\mathrm{ch}E_d}
   =
   \bigl(
   \operatorname{Ob}^{\mathrm{ch}\,E_2}_{\mathrm{Deligne}},
   \mathrm{Ob}_{F},
   \mathrm{Ob}_{R}
   \bigr),
\]
where \(\mathrm{Ob}_{F}\) and \(\mathrm{Ob}_{R}\) measure the obstruction
to the formality and Ran-descent hypotheses beyond the formal
Hamiltonian-current suboperad.  The formal-Darboux stalk theorem and
Theorem~\ref{app:thm:formal-polydisk-all-arity-current-operad}
compute the \(d=2\) local Hamiltonian-current operations on which any
global theorem restricts; vanishing of
\(\mathrm{Ob}_{\mathrm{ch}E_d}\) is the additional global condition.
\end{rmk}

\subsection{Descent over the Costello--Gwilliam brane-link boundary}\label{app:hfc:CG-descent}

The boundary open category $\mathcal{C}^{\mathrm{op}}_\partial$
(Definition~\ref{app:defn:brane-cat}) is defined globally on the
Costello--Gwilliam brane-link boundary $\partial X$.  This subsection
records the Weiss descent datum for local boundary categories.  A
global mixed-HT Deligne equivalence requires this descent datum
together with the chiral \(E_d\), quantum master equation, and
comparison-cone data of Definition~\ref{defn:oca-comparison-cone}.

\begin{defn}[Costello--Gwilliam brane-link bundle]\label{app:defn:brane-link}
For $X=\R^2_{\mathrm{top}}\times Y_{\mathrm{hol}}$ a mixed-HT
spacetime and $L\subset X$ a Dirac brane defect, the
\emph{Costello--Gwilliam brane-link boundary} $\partial X$ is the
boundary of the excised tubular neighbourhood $\nu(L)$ of $L$ in $X$
\cite{costello-gwilliam}*{Vol.~II},
\[
  \partial X
   \;=\;\partial(X\setminus\nu(L))
   \;\cong\;
   L\times S^{c-1}_\epsilon,
\]
the unit sphere bundle of codimension $c=\mathrm{codim}_X L$ over $L$.
At $L=\R_t\times\{0_s\}\times\{p\}
\subset\R^2_{\mathrm{top}}\times\C^2_{\mathrm{hol}}$
the codimension is $c=5$ and
$\partial X\cong\R_t\times S^4_\epsilon$.
\end{defn}

\begin{defn}[Collared formal-Darboux Weiss cover]\label{app:defn:collared-darboux-weiss-cover}
A \emph{Weiss cover} \(\mathfrak U=\{U_a\}_{a\in A}\) of
\(\partial X\) is an open cover such that every finite subset of
\(\partial X\) lies in some \(U_a\).  It is a
\emph{collared formal-Darboux Weiss cover} if each \(U_a\) is contained
in the link of a tubular chart
\[
  \nu(L)|_{I_a}\cong I_a\times
  \bigl(\R_s\oplus\C_{z_1}\oplus\C_{z_2}\bigr)
\]
over an interval \(I_a\subset L\), with the holomorphic normal directions
identified by a formal Darboux coordinate system and with the boundary
collar included.  On a multiple intersection
\(U_{a_0\cdots a_p}=U_{a_0}\cap\cdots\cap U_{a_p}\) the transition maps
are formal symplectomorphisms in the \((z_1,z_2)\)-directions and
ordinary collar-preserving smooth maps in the topological normal
direction.
\end{defn}

\begin{defn}[Boundary-category Weiss diagram on $\partial X$]\label{app:defn:sheaf-fact-cat-CG}
The boundary open category is locally encoded by a Weiss descent diagram
$\underline{\mathcal{C}^{\mathrm{op}}_\partial}$ of mixed-HT
factorization categories on $\partial X$: for
$U\subset\partial X$ open,
\[
  \underline{\mathcal{C}^{\mathrm{op}}_\partial}(U)
   =\mathcal C^{\mathrm{op}}_\partial(U),
\]
the category of boundary conditions supported on \(U\), together with
the restricted bulk--boundary action.  A module-category model over
\(\mathcal A^{\mathrm{cl}}_{\mathrm{bulk}}|_U\) is an additional
presentation.  At a closed point $p\in\partial X$ the formal fibre is
the formal-Darboux stalk of
Definition~\ref{app:defn:Darboux-stalk-functor}.
\end{defn}

\begin{thm}[Weiss descent criterion for the boundary open
category]\label{app:thm:cg-etale-descent}
Assume the boundary category is presented by a Costello--Gwilliam
boundary factorization category and satisfies the Weiss cosheaf
condition.  Then, for any Weiss cover
\(\{U_a\subset\partial X\}\), the homotopy-colimit descent map
\[
   \operatornamewithlimits{hocolim}\bigl(\textstyle\coprod_a
    \underline{\mathcal{C}^{\mathrm{op}}_\partial}(U_a)
    \rightrightarrows
    \coprod_{a,b}\underline{\mathcal{C}^{\mathrm{op}}_\partial}(U_a\cap U_b)
    \rightrightarrows\cdots\bigr)
  \;\xrightarrow{\sim}\;
  \underline{\mathcal{C}^{\mathrm{op}}_\partial}(\partial X)
\]
is an equivalence.
\end{thm}

\begin{proof}
This is the Weiss cosheaf axiom for Costello--Gwilliam
factorization categories on a real manifold, applied to the
brane-link boundary.  The mixed-HT labels and the bulk--boundary
action restrict to opens; compatibility with the displayed
homotopy-colimit map is part of the descent datum.
\end{proof}

\begin{defn}[Čech--Weiss descent obstruction complex]\label{app:defn:descent-obstruction-complex}
Choose a collared formal-Darboux Weiss cover
\(\mathfrak U=\{U_a\subset\partial X\}_{a\in A}\).  Put
\[
  Z_\partial=
  Z^{\mathrm{der},\mathrm{Mor}}_{E_2^{\mathrm{HT}}}
  (\mathcal C^{\mathrm{op}}_\partial),
  \qquad
  \mathcal A=\mathcal A^{\mathrm{cl}}_{\mathrm{bulk}}.
\]
For every non-empty multiple intersection
\(U_{a_0\cdots a_p}\), define the local mapping complex
\[
  \mathcal E^\bullet_{\mathrm{desc}}(U_{a_0\cdots a_p})
  =
  \underline{\mathrm{Hom}}^\bullet_{\mathcal B^{\mathrm{adm}}}
  \bigl(\mathcal A|_{U_{a_0\cdots a_p}},
        Z_\partial|_{U_{a_0\cdots a_p}}\bigr),
\]
with differential
\[
  d_{\mathcal E}\varphi
  =
  d_{Z_\partial}\varphi-(-1)^{|\varphi|}\varphi d_{\mathcal A}.
\]
The Čech--Weiss descent complex is the total complex
\[
  \check C^p_{\mathrm{Weiss}}(\mathfrak U;
  \mathcal E^q_{\mathrm{desc}})
  =
  \prod_{a_0<\cdots<a_p}
  \mathcal E^q_{\mathrm{desc}}(U_{a_0\cdots a_p}),
  \qquad
  d_{\mathrm{tot}}=\delta_{\check C}+(-1)^p d_{\mathcal E}.
\]

Let \(F_a\in\mathcal E^0_{\mathrm{desc}}(U_a)\) be the local
formal-Darboux closed-to-open maps obtained from the trace-sector stalk
and the local Morita \(E_2\)-centre lift.  A first overlap gluing datum
is a Čech \(1\)-cochain
\[
  H_{ab}\in \mathcal E^{-1}_{\mathrm{desc}}(U_{ab}),
  \qquad
  d_{\mathcal E}H_{ab}=F_b|_{U_{ab}}-F_a|_{U_{ab}}.
\]
The first descent obstruction is the triple-overlap class
\[
  \operatorname{Ob}_{\mathrm{desc}}
  =
  [\delta_{\check C}H]
  \in
  \check H^2_{\mathrm{Weiss}}
  \bigl(\mathfrak U;H^{-1}(\mathcal E^\bullet_{\mathrm{desc}})\bigr),
  \qquad
  (\delta_{\check C}H)_{abc}=H_{bc}-H_{ac}+H_{ab}.
\]
Equivalently, it is the edge class of \(\delta_{\check C}H\) in the
spectral sequence of
\(\operatorname{Tot}\check C^\bullet_{\mathrm{Weiss}}
(\mathfrak U;\mathcal E^\bullet_{\mathrm{desc}})\).  Higher descent
obstructions are the successive total-degree obstructions to extending
\((F_a,H_{ab})\) to a full homotopy-coherent Čech--Weiss descent datum.
\end{defn}

\begin{thm}[Descent obstruction at the formal-Darboux strata]\label{app:thm:global-stalk-from-collapse}
If the class \(\operatorname{Ob}_{\mathrm{desc}}\) of
Definition~\ref{app:defn:descent-obstruction-complex} and all higher
classes in
\(\operatorname{Tot}\check C^\bullet_{\mathrm{Weiss}}
(\mathfrak U;\mathcal E^\bullet_{\mathrm{desc}})\) are represented and
killed by compatible cochains, then the local trace-sector
formal-Darboux maps glue to a global morphism
\[
  \mathrm{OCA}_N\colon
   \mathcal A^{\mathrm{cl}}_{\mathrm{bulk}}
   \longrightarrow
   Z^{\mathrm{der},\mathrm{Mor}}_{E_2^{\mathrm{HT}}}
   \bigl(\mathcal{C}^{\mathrm{op}}_\partial\bigr)
\]
on the Costello--Gwilliam brane-link boundary.  It is a global
open--closed quasi-isomorphism only after the remaining chiral \(E_d\),
quantum master equation, and comparison-cone obstruction of
Definition~\ref{defn:oca-comparison-cone} vanish.
\end{thm}

\begin{proof}
Theorem~\ref{app:thm:stalk-of-conj} gives the trace-sector map on
each \(U_a\).  On double overlaps the chosen comparison homotopies are
the cochain \(H_{ab}\).  The equation
\(d_{\mathcal E}H_{ab}=F_b-F_a\) makes \(\delta_{\check C}H\) a
\(d_{\mathcal E}\)-cycle on triple overlaps.  Its class is independent
of changing \(H\) by a Čech coboundary or by a
\(d_{\mathcal E}\)-exact homotopy.  A null-homotopy of this class
chooses the next coherent gluing cochain; the same argument in the
total Čech--Weiss complex gives the higher obstruction tower.  Vanishing
of the tower is precisely the homotopy-coherent descent datum for the
local maps \(F_a\), hence produces the global morphism.  Its being a
quasi-isomorphism is the separate comparison-cone condition of
Definition~\ref{defn:oca-comparison-cone}.
\end{proof}

\subsection{All-order QME obstruction criterion for finite-jet PVA}\label{app:hfc:all-order-qme}

The classical PVA Jacobi identity holds at the formal-Darboux stalk
(Theorem~\ref{thm:reduced-ce-pv-central-operation}) and is the Jacobi
identity for the native two-dimensional singular product of
Theorem~\ref{thm:formal-two-dimensional-holomorphic-ope}.  Its quantum
lift is a recursive quantum master equation in the Costello
local-functional complex.  At order \(\hbar^n\) the obstruction lies in
the cohomology of the quantum BV differential with scalar-contact
projection, curved bulk-to-defect kernel, and compatible counterterms;
it is not killed by the classical Jacobi identity.  The hypotheses below
name the exact additional data needed to transport this recursion to the
trace-sector centre.

\paragraph{Convention.}
Under the binding convention of
\nameref{ssec:local-theorem-conventions}, the symbol
$C^\bullet_{\mathrm{ch}}(A_b,A_b)$ in this subsection denotes the
brane-side Hochschild target receiving the trace-sector map
\(\Phi_N^{\mathrm{gen}}\) of Theorem~\ref{thm:main-local}.  A full identification
with the bulk stalk is the admissible centre-lift problem.
\(Q_{\mathrm{ch}}\) is the Hochschild differential on this target.  It
is transported to the trace summand only after the admissible
Hochschild-centre lift, comparison-cone contraction, and required
homotopy-inverse data have been supplied.  The hypotheses below are
formulated on that target.

\begin{defn}[Hypothesis K: filtered formal SDR]\label{app:hyp:KZ-SDR}
There is a pronilpotent filtered strong deformation retract
\[
  \bigl(C^\bullet_{\mathrm{CE},\mathrm{cont}}(\mathfrak g),Q_{\mathrm{tot}}\bigr)
   \xrightarrow[\text{retract}]{\text{deform}}
   \bigl(C^\bullet_{\mathrm{ch}}(A_b,A_b),Q_{\mathrm{ch}}\bigr)
\]
over \(\C[[\hbar]]\), compatible with the scalar-contact filtration,
finite PVA weight windows, and the admissible Hochschild centrality
homotopies.  This is additional homological datum, not a consequence of
the formal-Darboux stalk theorem.
\end{defn}

\begin{rmk}[Formal and analytic transfer]\label{app:rmk:why-analytic}
The classical PVA Jacobi holds in the $\hbar$-adic topology of the
formal disk.  Costello's renormalization theorem gives, after a
renormalization scheme has been fixed, a finite-order obstruction
calculus.  A formal SDR transfers the obstruction classes.  It does not
prove analytic convergence.  Sectorial analytic convergence is the extra
content of Hypothesis~\ref{app:hyp:stokes} together with explicit graph
estimates.
\end{rmk}

\begin{defn}[Hypothesis S: sectorial Stokes regularity]\label{app:hyp:stokes}
If an analytic refinement is sought, the renormalized graph sums and the
transferred homological perturbation data form a sheaf over sectors
\(0<|\hbar|<r\).  On the sectors meeting the real-positive
\(\hbar\)-axis the Stokes automorphisms act trivially on the chosen
trace-sector image, and the sectorial lifts have the same asymptotic
\(\hbar\)-adic series.  No KZ connection is assumed unless it is
constructed as part of this analytic sheaf.
\end{defn}

\begin{defn}[Hypothesis W: finite-window weight symmetry]\label{app:hyp:reflected-weights}
Under crossing $z\leftrightarrow w$ on the formal \(E_2\) singular product of
Theorem~\ref{thm:formal-two-dimensional-holomorphic-ope}, the PVA
weight grading on $C^\bullet_{\mathrm{ch}}(A_b,A_b)$ reflects: the
operator
\[
  R\colon\mathrm{wt}_z\otimes\mathrm{wt}_w
   \longrightarrow\mathrm{wt}_w\otimes\mathrm{wt}_z
\]
intertwines the arity-two singular product with itself in each finite
weight window.  This is \(S_2\)-equivariance of the binary kernel; it is
not reflection positivity and it does not construct the all-arity
chiral \(E_2\) action.
\end{defn}

\begin{defn}[Hypothesis G: \texorpdfstring{$T=[Q_{\mathrm{tot}},G]$}{T = [Q,G]}]\label{app:hyp:T-Q-G}
In the Costello local-functional complex the holomorphic stress tensor is
quantum-BV exact modulo contact terms:
\[
  T=d_{\mathrm{QME}}^{L}G .
\]
The primitive \(G\) is a degree-\((-1)\) local functional compatible with
the same filtration and scalar-contact projection used in the QME
complex.  Transport to Hochschild cochains is allowed only after the
admissible quantum closed-to-open map has been constructed.
\end{defn}

\begin{thm}[All-order QME obstruction criterion for the trace-sector assignment]\label{app:thm:relative-qme}
Assume Hypotheses~\ref{app:hyp:KZ-SDR},
\ref{app:hyp:reflected-weights}, and \ref{app:hyp:T-Q-G}.  Assume also
the brane-defect Costello local-functional complex with
\(d_{\mathrm{QME}}^{L}\)-stable scalar-contact projection, finite-window
residual complexes, curved bulk-to-defect kernel, compatible local
counterterms, zero Milnor/Roos obstruction, and quantum centrality
homotopies for the centre lift.  If the Costello obstruction classes vanish in every loop
order with compatible primitives, the
formal-Darboux stalk of Theorem~\ref{app:thm:stalk-of-conj} admits a
formal all-order quantization: there exists a quantum BV master action
$S^{\mathrm{quant}}=S^{\mathrm{cl}}+\sum_{n\geq 1}\hbar^n S^{(n)}$
satisfying
\[
  QS^{\mathrm{quant}}
  +
  \frac12\{S^{\mathrm{quant}},S^{\mathrm{quant}}\}_L
  +\hbar\Delta_L S^{\mathrm{quant}}=0
\]
to all orders in \(\hbar\), and a quantum trace-sector closed-to-open
factorization-centre assignment
			\[
			  \Phi_N^\hbar\colon
			  \mathcal A^{\mathrm{q,tr}}_{\mathrm{bulk}}|_p
			   \longrightarrow
			   Z^{P_{0,\hbar}}_{\mathrm{fact}}
			   (\mathrm{Obs}^\hbar_{\mathrm{open}})|_b
				\]
					whose classical limit is the chosen lift
				\(\Phi_N^{\mathrm{Hoch}}\).  After a
quantum boundary \(A_\infty\)-algebra
\(A^\hbar_{\partial,N}\) has been chosen, this target may be modelled by
\[
  C^\bullet_{\mathrm{Hoch},\mathrm{adm}}
  (A^\hbar_{\partial,N},A^\hbar_{\partial,N}).
\]
The trace
map \(J(f)=\overline{\operatorname{Tr}\,f(\phi_1,\phi_2)}\) is defined
on polynomials and extended to formal \(f\) after completion; its
		quantum correction has the form
		\(J^{\mathrm{quant}}(f)=J(f)+\sum_n\hbar^n J^{(n)}(f)\).
			The assignment is a morphism of Hochschild complexes only in the
chosen Hochschild model and precisely after the quantum
equivariance homotopies are chosen.  Descent to Hochschild cohomology and an
all-order derived-centre equivalence require the
			additional Hochschild equivariance, centre-lift, and descent data of
			Definition~\ref{defn:admissible-hochschild-centre-lift}.
			\end{thm}

\begin{proof}
Set \(S^{(0)}=S^{\mathrm{cl}}\).  At order \(n\), after
\(S^{(1)},\ldots,S^{(n-1)}\) have been chosen, the
order-\(n\) part of the quantum master equation has the form
\[
  d_{\mathrm{cl}}^L S^{(n)}
  =-\frac12\sum_{\substack{i+j=n\\ i,j<n}}\{S^{(i)},S^{(j)}\}_L
   -\Delta_L S^{(n-1)},\qquad
  d_{\mathrm{cl}}^L=Q+\{S^{(0)}[L],-\}_L .
\]
The right-hand side is \(d_{\mathrm{cl}}^L\)-closed by the Jacobi identity
for the BV bracket and by \(\Delta_L^2=0\).  Its cohomology class is the
Costello obstruction at loop order \(n\).  By hypothesis this class
vanishes and a compatible primitive is chosen; induction constructs
\(S^{(n)}\) for every \(n\).

The scalar-contact projection, finite-window residual complexes, curved
bulk-to-defect kernel, compatible counterterms, and zero Milnor/Roos
class supply the primitives in the brane-defect tower.
Hypothesis~\ref{app:hyp:KZ-SDR} transfers the formal tower to the
trace-sector centre.  Hypothesis~\ref{app:hyp:reflected-weights} is the
finite-window \(S_2\)-compatibility of the binary kernel.
Hypothesis~\ref{app:hyp:T-Q-G} supplies the local-functional homotopy for
the stress tensor.  Hypothesis~\ref{app:hyp:stokes}, when imposed, gives
sectorial analytic compatibility; it is not used to create a formal
primitive.
\end{proof}

\begin{rmk}[QME obstruction class]\label{app:rmk:qme-obstruction}
Before comparison maps to a common local-functional complex are fixed, the
QME obstruction classes form the tuple
\[
  \mathrm{Ob}_{\mathrm{QME}}
   =
   \bigl([\mathfrak{ob}_K],[\mathfrak{ob}_S],
   [\mathfrak{ob}_W],[\mathfrak{ob}_G]\bigr),
\]
with entries the formal-SDR obstruction \( [\mathfrak{ob}_K]\), the
sectorial Stokes class \( [\mathfrak{ob}_S]\), the binary
weight-symmetry class \( [\mathfrak{ob}_W]\), and the
stress-tensor-exactness class \( [\mathfrak{ob}_G]\).  It becomes a
single class in \(H^*(\mathrm{Loc}_{\R^2\times\C^2})\) only after these
entries have been represented in the same Costello local-functional
deformation complex.  At fixed loop order, only the algebraic finite-window
entries are checked in finite PVA weight windows.  Analytic convergence
requires the sectorial estimates of Hypothesis~\ref{app:hyp:stokes}.
\end{rmk}

\subsection{Operator-level lift of the Igusa section and its square-root branch}\label{app:hfc:operator-lift-modular}

The Igusa cusp form $\Phi_{10}$ is the level-2 section of the modular
line $\Omega_{\mathrm{central}}$ on the Siegel period domain
$\mathfrak H_2$
\cite{igusa-cusp-form}\cite{gritsenko-nikulin-paramodular}.  The square
root $\Delta_5=\Phi_{10}^{1/2}$ is a scalar modular form in the
compact \(K3\times E\) comparison
(\S\ref{ssec:open-modular-lift}).  The formal-Darboux theorem
constructs the local Capelli scalar.  An operator-level lift is
an operator-valued cocycle whose protected trace is \(\Phi_{10}\) on
the Siegel period domain, or \(\Delta_5\) after the square-root cover
has been chosen, in the compact-target theory.

\begin{conj}[Operator lift of the Igusa modular section]\label{app:conj:operator-lift-Delta5}
There exists a graded vertex-algebraic or factorization-algebraic
object \(V\), an involution or real form \(\sigma\), and a protected
trace such that
\[
  \mathrm{Tr}_V\bigl(\sigma\,q^{L_0-c/24}\,
   \prod_{i=1}^{2}z_i^{H_i}\bigr)
   \;=\;S(\tau,\rho,\nu),
   \qquad
   S\in\{\Phi_{10},\,\Delta_5=\Phi_{10}^{1/2}\}
\]
with \(S=\Delta_5\) only on the square-root cover of the Siegel period domain
$(\tau,\rho,\nu)\in\mathfrak H_2$, with the trace compatible with
the modular line \(\Omega_{\mathrm{central}}\).  The required object is
the operator-valued cocycle lifting this scalar modular section.
\end{conj}

Scalar automorphy does not determine an operator-valued cocycle.  The
missing datum is a vertex algebra or a factorization algebra with a
protected trace.

\begin{constr}[Borcherds--Kac--Moody lift obstruction]\label{app:constr:BKM-obstruction}
The Borcherds--Kac--Moody denominator formulas give automorphic
products, including the fake-Monster denominator and the
Gritsenko--Nikulin products
\cite{borcherds-fake}\cite{gritsenko-nikulin-paramodular}.  These are
scalar identities.  Passing from a scalar denominator formula to a
vertex-algebraic object whose graded trace is \(\Delta_5\) requires an
additional construction.

\emph{Obstruction.}
\(\mathrm{Ob}_{\mathrm{BKM}}\) is the obstruction to lifting the
scalar Borcherds product to an operator trace with the required
integral form, real structure, and modular covariance.
\end{constr}

\begin{constr}[Free-field realisation on
\texorpdfstring{$K3\times E$}{K3 times E}]\label{app:constr:K3E-obstruction}
The elliptic genus of \(K3\) \cite{witten-elliptic-genus}, its
second-quantised product, and the Donaldson--Thomas /
Pandharipande--Thomas comparison on
\(K3\times E\) \cite{maulik-pandharipande-thomas} give scalar modular data.
The obstruction measures the failure to construct a free-field or
factorization algebra on \(K3\times E\) whose protected trace is
\(\Delta_5\).

\emph{Obstruction.}
\(\mathrm{Ob}_{K3E}\) is the obstruction to constructing the compact
target factorization object, its protected trace, and the descent of
that trace to the level-\(2\) modular line.
\end{constr}

\begin{constr}[Dabholkar--Murthy--Zagier lift obstruction]\label{app:constr:DMZ-obstruction}
The Dabholkar--Murthy--Zagier decomposition of meromorphic Jacobi
forms separates polar and finite parts in the \(1/\Phi_{10}\)
dyonic counting problem \cite{dabholkar-murthy-zagier}.  This is a
statement about mock modular and wall-crossing structures.  An
operator construction of \(\Delta_5\) is the lift to the modular-line
cocycle.

\emph{Obstruction.}
\(\mathrm{Ob}_{\mathrm{DMZ}}\) is the obstruction to lifting the
mock-modular wall-crossing structure to an honest operator trace with
the required holomorphic modular transformation law.
\end{constr}

\begin{rmk}[Composition of operator-lift obstructions]\label{app:rmk:operator-lift-obstructions}
The scalar theories determine modular forms in the same comparison
class, and the operator-valued cocycle is the common lifting datum.  The symbol
\(\mathrm{Ob}_V\) denotes the direct sum of the named
obstruction complexes above.  A cohomological product formula for these
scalar obstruction classes requires a product on the direct-sum
obstruction complex.
\end{rmk}

\subsection{Obstruction taxonomy}\label{app:hfc:obstruction-taxonomy}

The local theorem and the obstruction classes assemble into the global
Mixed Holomorphic--Topological Deligne criterion.  The criterion is a
vanishing theorem for obstruction data; it is not a consequence of the
formal-Darboux trace-sector computation alone.

\begin{thm}[Global Mixed-HT Deligne criterion]\label{app:thm:global-deligne-criterion}
Let \(\mathrm{preOCA}_N\) be a pre-open--closed comparison datum whose
formal-Darboux fibre is the trace-sector morphism of
Theorem~\ref{app:thm:stalk-of-conj}.  Under the joint vanishing
\[
	  \mathrm{Ob}_{\mathrm{global}}
	   :=
	   \bigl(
	     \operatorname{Ob}^{\mathrm{ch}\,E_2}_{\mathrm{Deligne}},
	     \mathrm{Ob}_{\mathrm{desc}},
	     \mathrm{Ob}_{\mathrm{QME}}
	   \bigr)
   =0,
\]
with the three summands as in
Definition~\ref{app:defn:morita-chiral-E2-extension-complex},
Theorem~\ref{app:thm:global-stalk-from-collapse},
and Remark~\ref{app:rmk:qme-obstruction}, and after the anomaly,
locality, and comparison-cone data of
Definition~\ref{defn:global-mixed-ht-deligne-obstruction-datum} are
killed by compatible null-homotopies, \(\mathrm{preOCA}_N\) is promoted
to a global open--closed quasi-isomorphism.  On this vanishing locus
with the displayed primitives chosen, this is the \(d=2\) realization of the
Mixed-HT Deligne conjecture~\ref{app:conj:mixed-HT-Deligne} on
$\R^2_{\mathrm{top}}\times\C^2_{\mathrm{hol}}$ at $d=2$, $N$ Dirac
branes, with every construction hypothesis named in those theorems.
The operator-level lift of \(\Phi_{10}\) or
\(\Phi_{10}^{1/2}\) belongs to the \(K3\times E\) comparison problem.
\end{thm}

\begin{proof}
Compose Proposition~\ref{app:prop:morita-centre-lift-MC-control}
(Morita chiral \(E_2\)-centre lift via the relative Maurer--Cartan
complex),
Theorem~\ref{app:thm:relative-Ed-deligne} (chiral $E_d$
Deligne via the operadic action and
Hypotheses~\ref{app:hyp:formality}--\ref{app:hyp:ran-descent}),
Theorem~\ref{app:thm:global-stalk-from-collapse}
(Costello--Gwilliam brane-link descent via
Definition~\ref{app:defn:descent-obstruction-complex}),
Theorem~\ref{app:thm:relative-qme} (all-order QME obstruction criterion via
Hypotheses~\ref{app:hyp:KZ-SDR}, \ref{app:hyp:reflected-weights}, and
\ref{app:hyp:T-Q-G}, with Hypothesis~\ref{app:hyp:stokes} only for
analytic refinement).  The three obstructions are independent cohomology
classes living in distinct cohomology groups; their tuple becomes a
single curvature class only after a common deformation complex and
comparison maps have been constructed.
\end{proof}

\begin{rmk}[Reduction to obstruction classes]\label{app:rmk:reduction-status}
The obstruction $\mathrm{Ob}_{\mathrm{QME}}$ is the all-order quantum
master-equation obstruction for lifting the finite-jet PVA across all
$\hbar$-orders.  The Morita \(E_2\)-Deligne obstruction
\(\operatorname{Ob}^{\mathrm{ch}\,E_2}_{\mathrm{Deligne}}\) is the class
of Definition~\ref{app:defn:morita-chiral-E2-extension-complex}; at
\(d\geq2\) it is the Dirac-brane formal-Darboux instance of the
obstruction class of Conjecture~\ref{conj:chiral-Ed-deligne};
  partial results (chiral Lie algebroids on smooth schemes via Hennion,
  factorisable co-operads via Francis) give adjacent structures; the
  global assertion is the obstruction criterion of
  Theorem~\ref{app:thm:global-deligne-criterion}, and the
  underlying equivalence remains
  Conjecture~\ref{app:conj:mixed-HT-Deligne} until that criterion is
  satisfied.  The
  brane-link descent obstruction $\mathrm{Ob}_{\mathrm{desc}}$ is the
  Čech class of Definition~\ref{app:defn:descent-obstruction-complex}.
	The formal-Darboux trace theorem
	(Theorem~\ref{thm:main-local}) is the formal-Darboux stalk
that anchors the local side of these obstructions, and
Theorem~\ref{thm:formal-two-dimensional-holomorphic-ope} gives its
arity-two formal diagonal operation.  Each relative theorem has the
stalk theorem, including this diagonal singular product, as its base case.
\end{rmk}

\begin{constr}[Global obstruction class]\label{app:constr:global-obstruction-class}
The global criterion
(Theorem~\ref{app:thm:global-deligne-criterion}) translates the
compact-comparison obstructions of
\S\ref{sec:compact-comparison-obstructions} into
cohomology classes in distinct source complexes:

\begin{itemize}
  \item
    \(\operatorname{Ob}^{\mathrm{ch}\,E_2}_{\mathrm{Deligne}}\):
    Maurer--Cartan lift from the formal tree-kernel current operad to
    the Morita mixed-HT chiral \(E_2\)-centre action
    (Definition~\ref{app:defn:morita-chiral-E2-extension-complex}); the
    higher \(d\geq2\) formality summand is the associator-type operadic
    homotopy class on \(\mathrm{Conf}_I(\C^d)\), \(|I|\geq2\).
  \item $\mathrm{Ob}_{\mathrm{desc}}$: null-homotopy of the Čech class
    in the descent obstruction complex
    (Definition~\ref{app:defn:descent-obstruction-complex}) on the
    formal-Darboux stratum.
  \item $\mathrm{Ob}_{\mathrm{QME}}$: each of four checkable
    cohomology classes
    (Hypotheses~\ref{app:hyp:KZ-SDR}--\ref{app:hyp:T-Q-G}).
\end{itemize}

The formal-Darboux stalk theorem
(Theorem~\ref{thm:main-local}) is the local trace-sector fibre.
The operator lift of the modular line on \(K3\times E\) is governed by
the compact-target comparison complex, rather than by a summand of
\(\mathrm{Ob}_{\mathrm{global}}\).
\end{constr}
